\section{Jacobson preschemes}
\label{section:IV.10}


\subsection{Very dense subsets of a topological space}
\label{subsection:IV.10.1}

\oldpage[IV-3]{95}

\begin{env}[10.1.1]
\label{IV.10.1.1}
We say that a subset of a topological space $X$ is \emph{quasi-constructible} if it is a finite union of locally closed subsets of $X$.
We say that a subset $T$ of $X$ is \emph{locally quasi-constructible} if for every $x\in X$, there exists an open neighbourhood $V$ of $x$ such that $T\cap V$ is quasi-constructible in $V$.
It is clear that every quasi-constructible subset of $X$ is locally quasi-constructible; the converse is true if $X$ is quasi-compact.
The argument of \sref[III]{III.9.1.3}, where one removes the word ``retrocompact'', shows that the set of quasi-constructible (resp.\ locally quasi-constructible) subsets of $X$, which we will denote by $\mathfrak{Qc}(X)$ (resp.\ $\mathfrak{Lqc}(X)$), is stable under finite union, finite intersection, and passage to complements.
If $f:X\to Y$ is a continuous map, it follows immediately from these definitions that, for every quasi-constructible (resp.\ locally quasi-constructible) subset $Z$ of $Y$, $f^{-1}(Z)$ is quasi-constructible (resp.\ locally quasi-constructible) in $X$.

The constructible (resp.\ locally constructible) subsets of $X$ are evidently quasi-constructible (resp.\ locally quasi-constructible); the converse is true when $X$ is Noetherian (resp.\ locally Noetherian).

In what follows, we will denote respectively by $\mathfrak{O}(X)$, $\mathfrak{F}(X)$, $\mathfrak{Lf}(X)$, $\mathfrak{C}(X)$, $\mathfrak{Lc}(X)$ the set of subsets of $X$ that are respectively open, closed, locally closed, constructible, locally constructible.
\end{env}

\oldpage[IV-3]{96}

\begin{proposition}[10.1.2]
\label{IV.10.1.2}
Let $X$ be a topological space, $X_0$ a subset of $X$.
The following conditions are equivalent:
\begin{enumerate}[label=\emph{\alph*)}]
  \item For every locally closed subset $Z\neq\varnothing$ of $X$, we have $Z\cap X_0\neq\varnothing$.
  \item[\rm a$'$)] For every closed subset $Z$ of $X$, we have $Z=\overline{Z\cap X_0}$.
  \item For every locally quasi-constructible subset $Z\neq\varnothing$ of $X$, we have $Z\cap X_0\neq\varnothing$.
  \item[\rm b$'$)] For every locally quasi-constructible subset $Z$ of $X$, we have $Z\subset\overline{Z\cap X_0}$ (in other words, $Z\cap X_0$ is dense in $Z$).
  \item The map $U\mapsto U\cap X_0$ of $\mathfrak{O}(X)$ into $\mathfrak{O}(X_0)$ is injective (hence bijective).
  \item[\rm c$'$)] The map $F\mapsto F\cap X_0$ of $\mathfrak{F}(X)$ into $\mathfrak{F}(X_0)$ is injective (hence bijective).
  \item[\rm c$''$)] The map $Z\mapsto Z\cap X_0$ of $\mathfrak{Qc}(X)$ into $\mathfrak{Qc}(X_0)$ is injective (hence bijective).
  \item[\rm c$'''$)] The map $Z\mapsto Z\cap X_0$ of $\mathfrak{Lqc}(X)$ into $\mathfrak{Lqc}(X_0)$ is injective.
\end{enumerate}
Furthermore, when these conditions are satisfied, the map $Z\mapsto Z\cap X_0$ of $\mathfrak{Lqc}(X)$ into $\mathfrak{Lqc}(X_0)$ is bijective.
\end{proposition}

\begin{proof}
Let us note that the assertions of surjectivity of $c)$ and $c')$ are trivial; they mean that every locally closed subset of $X_0$ is the trace on $X_0$ of a locally closed subset of $X$, and that the map $\mathfrak{Qc}(X)\to\mathfrak{Qc}(X_0)$ defined in $c'')$ is also surjective.

We will prove the implications
\[
  c''')\Rightarrow c'')\Rightarrow c')\Rightarrow c)\Rightarrow b)\Rightarrow b')\Rightarrow a')\Rightarrow a)\Rightarrow c''').
\]

The first three are trivial.
To see that $c)$ implies $b)$, let us first note that $c)$ means that $X_0$ is \emph{dense} in $X$.
Replacing $X$ and $X_0$ by $U$ and $U\cap X_0$ respectively, one can therefore restrict to the case where $Z$ is locally closed in $X$, hence $Z=V\cap\complement W$, where $V$ and $W$ are open in $X$; the hypothesis $Z\neq\varnothing$ means that $V\not\subset W$, or equivalently $V\cup W\neq W$.
By virtue of $c)$, we have $(V\cup W)\cap X_0\neq W\cap X_0$, hence $V\cap X_0\not\subset W$, and consequently $(V\cap\complement W)\cap X_0\neq\varnothing$.

To see that $b)$ implies $b')$, it suffices to apply $b)$ to $Z\cap U$, where $U$ is an open neighbourhood of an arbitrary point of $Z$.
As $Z\cap\overline{X_0}\subset\overline{Z}$, it is trivial that $b')$ implies $a')$.
To show that $a')$ implies $a)$, let us note that if $Z$ is locally closed in $X$, we can write $Z=F-F'$ where $F$ and $F'$ are closed in $X$ and $F'\subset F$; whence $Z\cap X_0=(F\cap X_0)-(F'\cap X_0)$.
If $Z\cap X_0=\varnothing$, then, by virtue of $a')$, we would have $F\cap X_0=F'\cap X_0$, hence $F=F'$, i.e.\ $Z=\varnothing$.

To see that $a)$ implies $c''')$, it suffices to show that if $Z\cap X_0\neq\varnothing$, then $Z'\cap X_0=Z''\cap X_0$ is equivalent to $((Z'\cup Z'')-(Z'\cap Z''))\cap X_0=\varnothing$.
In other words, it suffices to prove that $a)$ implies $b)$; indeed, replacing $X$ and $X_0$ by $U$ and $U\cap X_0$ respectively, one is reduced to the case where $Z$ is locally closed in $X$, whence the conclusion.

It remains to show that the map $\mathfrak{Lqc}(X)\to\mathfrak{Lqc}(X_0)$ is surjective.
So let $Z_0$ be a locally quasi-constructible subset of $X_0$: there is a covering $(V_\alpha)$ of $X_0$ by open subsets of $X_0$ such that $Z_0\cap V_\alpha$ is quasi-constructible in $V_\alpha$ (and
\oldpage[IV-3]{97}
hence also in $X_0$).
By virtue of $c)$, there is for each $\alpha$ a unique quasi-constructible subset $Z_\alpha$ in $X$ such that $Z_0\cap V_\alpha=X_0\cap Z_\alpha$.
If $\alpha$ and $\beta$ are any two indices, we have $Z_\beta\cap U_\alpha\cap X_0=U_\alpha\cap X_0=V_\alpha\cap Z_0\cap V_\beta=Z_\alpha\cap X_0\cap V_\beta\subset Z_0\cap X_0$; as $Z_\beta\cap U_\alpha$ and $Z_\alpha$ are quasi-constructible in $X$, it follows from $c'')$ that $Z_\beta\cap U_\alpha\subset Z_\alpha$.
If we set $Z=\bigcup_\alpha Z_\alpha$, then $Z\cap U_\alpha=Z_\alpha$ for every $\alpha$; moreover, since the $V_\alpha$ cover $X_0$, it follows from $c)$ that the $U_\alpha$ cover $X$, and we see therefore that $Z$ is locally quasi-constructible in $X$ and $Z_0=Z\cap X_0$.
\end{proof}

\begin{definition}[10.1.3]
\label{IV.10.1.3}
When a subset $X_0$ of a topological space satisfies the equivalent conditions of \sref{IV.10.1.2}, we say that $X_0$ is \emph{very dense} in $X$.
\end{definition}

We have already seen in the course of the proof of \sref{IV.10.1.2} that $X_0$ is then dense in $X$.

\begin{corollary}[10.1.4]
\label{IV.10.1.4}
If $X_0$ is very dense in $X$ and $U$ an open subset of $X$, $U\cap X_0$ is very dense in $U$.
Conversely, if $(U_\alpha)$ is an open covering of $X$ such that $U_\alpha\cap X_0$ is very dense in $U_\alpha$ for every $\alpha$, then $X_0$ is very dense in $X$.
\end{corollary}

As every locally closed subset of $U$ is locally closed in $X$, the first assertion follows from criterion $a)$ of \sref{IV.10.1.2}; and the same is true of the second, since if $Z\neq\varnothing$ is locally closed in $X$, $Z\cap U_\alpha$ is locally closed in $U_\alpha$ for every $\alpha$, and $Z\cap U_\alpha\neq\varnothing$ for at least one $\alpha$.

\subsection{Quasi-homeomorphisms}
\label{subsection:IV.10.2}

\begin{proposition}[10.2.1]
\label{IV.10.2.1}
Let $X_0$, $X$ be two topological spaces, $f:X_0\to X$ a continuous map.
The following conditions are equivalent:
\begin{enumerate}[label=\emph{\alph*)}]
  \item The map $U\mapsto f^{-1}(U)$ of $\mathfrak{O}(X)$ into $\mathfrak{O}(X_0)$ is bijective.
  \item[\rm a$'$)] The map $F\mapsto f^{-1}(F)$ of $\mathfrak{F}(X)$ into $\mathfrak{F}(X_0)$ is bijective.
  \item The topology of $X_0$ is the inverse image by $f$ of that of $X$, and $f(X_0)$ is very dense \sref{IV.10.1.3} in $X$.
  \item The functor $\sh{F}\mapsto f^*(\sh{F})$ from the category of sheaves of sets (resp.\ sheaves of abelian groups) on $X$ to the category of sheaves of sets (resp.\ sheaves of abelian groups) on $X_0$ is an equivalence of categories.
\end{enumerate}
\end{proposition}

\begin{proof}
It is clear that $a)$ and $a')$ are equivalent, and $a)$ implies that the topology of $X_0$ is the inverse image by $f$ of that of $X$.
On the other hand, if $f(X_0)$ is not very dense in $X$, there are two distinct open subsets $U_1$, $U_2$ of $X$ such that $U_1\cap f(X_0)=U_2\cap f(X_0)$, and hence $f^{-1}(U_1)=f^{-1}(U_2)$, which completes showing that $a)$ implies $b)$.
Conversely, the condition $b)$ implies that the maps $U\mapsto U\cap f(X_0)$ of $\mathfrak{O}(X)$ into $\mathfrak{O}(f(X_0))$ and $V\mapsto f^{-1}(V)$ of $\mathfrak{O}(f(X_0))$ into $\mathfrak{O}(X_0)$ are bijective, hence so is the composite $U\mapsto f^{-1}(U)$.

To see that $a)$ implies $c)$, it suffices to apply the definition of $f^*(\sh{F})$ \sref[0]{0.3.7.1} and the axioms for sheaves: $a)$ means that for every open $U$ of $X$, the canonical map $\Gamma(U,\sh{F})\to\Gamma(f^{-1}(U),f^*(\sh{F}))$ is a functorial bijection in $\sh{F}$, whence $c)$.
It remains to show that $c)$ implies $b)$.

\oldpage[IV-3]{98}
Suppose first that $f(X_0)$ is not very dense in $X$; there then exist two distinct closed subsets $Y\supset Y'$ of $X$ such that $Y\cap f(X_0)=Y'\cap f(X_0)$.
Let $\sh{F}$ (resp.\ $\sh{F}'$) be the sheaf of abelian groups on $X$ that is the direct image by the canonical injection $Y\to X$ (resp.\ $Y'\to X$) of the simple sheaf associated to the constant presheaf $\mathbf{Z}$ on $Y$ (resp.\ $Y'$); the definition of $f^*$ \sref[0]{0.3.7.1} shows that $f^*(\sh{F})$ is isomorphic to $f^*(\sh{F}')$, but $\sh{F}$ is not isomorphic to $\sh{F}'$, hence condition $c)$ is not satisfied.
(One notes that the functor $f^*$ is not then even \emph{faithful}, for if $u$ and $v$ are the identity automorphism and the zero endomorphism of $\sh{F}/\sh{F}'$, $f^*(u)$ and $f^*(v)$ are equal.)

Let us now show that $c)$ implies that the topology of $X_0$ is the inverse image of that of $X$ by $f$.
Let us note that if condition $c)$ is satisfied for the category of sheaves of sets, it is also satisfied for the category of sheaves of abelian groups, for it follows immediately from the definitions \sref[III]{III.8.2.5} that the category of \emph{objects in abelian groups} in the category of sheaves of sets is none other than the category of sheaves of abelian groups.
It will therefore suffice to prove our assertion supposing $c)$ satisfied for the categories of sheaves of abelian groups on $X$ and $X_0$.
Now, if $\mathbf{Z}_X$ denotes the simple sheaf on $X$ associated to the constant presheaf $\mathbf{Z}$, there is an isomorphism canonical up to sign $\Gamma(X,\sh{F})\xrightarrow{\sim}\operatorname{Hom}(\mathbf{Z}_X,\sh{F})$ functorial in $\sh{F}$ (by the same argument as in \sref[0]{0.5.1.1}); as it is clear by definition \sref[0]{0.3.7.1} that $f^*(\mathbf{Z}_X)=\mathbf{Z}_{X_0}$, the hypothesis that $f^*$ is an equivalence means that the canonical homomorphism $\Gamma(X,\sh{F})\to\Gamma(X_0,f^*(\sh{F}))$ is \emph{bijective} for every sheaf of abelian groups $\sh{F}$ on $X$.
Now, let $Y_0$ be a closed subset of $X_0$; let $\sh{F}_0$ be the sheaf on $X$ that is the direct image by the canonical injection $Y_0\to X_0$ of the simple sheaf associated to the constant presheaf $\mathbf{Z}$ on $Y_0$.
As $f^*$ is an equivalence, $\sh{F}_0$ is isomorphic to a sheaf of the form $f^*(\sh{F})$, where $\sh{F}$ is a sheaf of abelian groups on $X$.
Consider then a section $s_0$ of $\sh{F}_0$ over $X_0$ such that $(s_0)_{x_0}=1_{x_0}$ if $x_0\in Y_0$, $(s_0)_{x_0}=0_{x_0}$ if $x_0\notin Y_0$.
The preceding remarks show that there exists a section $s$ of $\sh{F}$ over $X$ such that $(s_0)_{x_0}=s_{f(x_0)}$ for every $x_0\in X_0$ (the fibres $(\sh{F}_0)_{x_0}$ and $\sh{F}_{f(x_0)}$ being canonically identified \sref[0]{0.3.7.1}); this means that $Y_0$ is the inverse image by $f$ of the set $Y$ of $x\in X$ such that $s_x\neq 0_x$, and $y$ is closed in $X$.
\end{proof}

\begin{definition}[10.2.2]
\label{IV.10.2.2}
When a continuous map $f:X_0\to X$ satisfies the equivalent conditions of \sref{IV.10.2.1}, we say that $f$ is a \emph{quasi-homeomorphism} of $X_0$ into $X$.
\end{definition}

By virtue of \sref{IV.10.2.1},~$b)$, to say that a subset $X_0$ of a topological space $X$ is very dense in $X$ means that the canonical injection $X_0\to X$ is a quasi-homeomorphism.

\begin{corollary}[10.2.3]
\label{IV.10.2.3}
The composite of two quasi-homeomorphisms is a quasi-homeomorphism.
\end{corollary}

This follows immediately from \sref{IV.10.2.1},~$a)$.

\begin{corollary}[10.2.4]
\label{IV.10.2.4}
Let $f:X\to Y$ be a quasi-homeomorphism, $Y'$ a locally quasi-constructible subset of $Y$, $X'=f^{-1}(Y')$; then the restriction $f'=f|X':X'\to Y'$ is a quasi-homeomorphism.
\end{corollary}

\begin{proof}
It is clear by virtue of \sref{IV.10.2.1},~$b)$ that the topology induced on $X'$ by that of $X$ is the inverse image by $f'$ of the topology induced on $Y'$ by that of $Y$.
On the other hand,
\oldpage[IV-3]{99}
let $Z\neq\varnothing$ be a closed subset \emph{in} $Y'$, and $z\in Z$; there is an open neighbourhood $U$ of $z$ in $Y$ such that $U\cap Y'$ is the union of a finite number of closed subsets $Y'_j$ of $U$; if $j$ is an index such that $z\in Y'_j$, $Z\cap Y'_j$ is therefore closed in $U$.
Since $f(X)\cap U$ is very dense in $U$ \sref{IV.10.1.4}, $Z\cap Y'_j\cap f(X)$ is not empty \sref{IV.10.1.2},~$a)$; but as $Z\subset Y'$, we have $Z\cap f(X)=Z\cap f'(X')$; the criterion \sref{IV.10.1.2},~$a)$ shows that $f'(X')$ is very dense in $Y'$, and we conclude by \sref{IV.10.2.1},~$b)$.
\end{proof}

\begin{corollary}[10.2.5]
\label{IV.10.2.5}
Let $f:X\to Y$ be a continuous map, $(V_\alpha)$ an open covering of $Y$.
If, for every $\alpha$, the restriction $f^{-1}(V_\alpha)\to V_\alpha$ of $f$ is a quasi-homeomorphism, then $f$ is a quasi-homeomorphism.
\end{corollary}

This follows immediately from the criterion \sref{IV.10.2.1},~$b)$ and from \sref{IV.10.1.4}.

\begin{corollary}[10.2.6]
\label{IV.10.2.6}
Let $f:X\to Y$ be a quasi-homeomorphism, $Y'$ a locally quasi-constructible subset of $Y$, $X'=f^{-1}(Y')$.
For $Y'$ to be quasi-compact (resp.\ Noetherian, resp.\ retrocompact in $Y$), it is necessary and sufficient that $X'$ be quasi-compact (resp.\ Noetherian, resp.\ retrocompact in $X$).
\end{corollary}

\begin{proof}
Let us prove the first two assertions; by virtue of \sref{IV.10.2.4}, one can restrict to the case where $Y'=Y$.
To say that $X$ is quasi-compact (resp.\ Noetherian) means that for every increasing filtered family $(U_\alpha)$ in $\mathfrak{O}(X)$ (resp.\ for every increasing filtered family $(U_\alpha)$ in $\mathfrak{O}(X)$) having $X$ as greatest element, there exists $\gamma$ such that $U_\alpha=U_\gamma$ for $\alpha\geq\gamma$.
As $U\mapsto f^{-1}(U)$ is a bijection of $\mathfrak{O}(Y)$ onto $\mathfrak{O}(X)$, our assertion follows immediately from the preceding remark.

The quasi-compact open subsets of $X$ are therefore the sets of the form $f^{-1}(U)$ where $U$ is an open quasi-compact subset of $Y$, by virtue of \sref{IV.10.2.1},~$a)$ and of what precedes.
For $X'$ to be retrocompact in $X$, it is necessary and sufficient that for every quasi-compact open $U$ in $Y$, $f^{-1}(U)\cap X'=f^{-1}(U\cap Y')$ be quasi-compact \sref[III]{III.9.1.1}; the first part of the proof shows that this amounts to saying that $U\cap Y'$ is quasi-compact for every open quasi-compact $U$, i.e.\ that $Y'$ is retrocompact in $Y$.
\end{proof}

\begin{proposition}[10.2.7]
\label{IV.10.2.7}
Let $f:X\to Y$ be a quasi-homeomorphism.
Then the map $Z\mapsto f^{-1}(Z)$ of $\mathfrak{P}(Y)$ into $\mathfrak{P}(X)$ defines by restriction the bijections (cf.\ \sref{IV.10.1.1} for the notation):
\[
  \mathfrak{O}(Y)\xrightarrow{\sim}\mathfrak{O}(X)
\]
\[
  \mathfrak{F}(Y)\xrightarrow{\sim}\mathfrak{F}(X)
\]
\[
  \mathfrak{Lf}(Y)\xrightarrow{\sim}\mathfrak{Lf}(X)
\]
\[
  \mathfrak{Qc}(Y)\xrightarrow{\sim}\mathfrak{Qc}(X)
\]
\[
  \mathfrak{Lqc}(Y)\xrightarrow{\sim}\mathfrak{Lqc}(X)
\]
\[
  \mathfrak{C}(Y)\xrightarrow{\sim}\mathfrak{C}(X)
\]
\[
  \mathfrak{Lc}(Y)\xrightarrow{\sim}\mathfrak{Lc}(X).
\]
\end{proposition}

\begin{proof}
\oldpage[IV-3]{100}
For the first two, this is none other than the definition \sref{IV.10.2.2}; as the topology of $X$ is the inverse image by $f$ of that of $f(X)$, the five remaining maps, where one replaces $Y$ by $f(X)$, are bijective.
It therefore suffices (by \sref{IV.10.2.1},~$b)$) to restrict to the case where $X$ is a very dense subspace of $Y$, and that the maps $\mathfrak{Lf}(Y)\to\mathfrak{Lf}(X)$, $\mathfrak{Qc}(Y)\to\mathfrak{Qc}(X)$, $\mathfrak{Lqc}(Y)\to\mathfrak{Lqc}(X)$ are bijective has already been proved \sref{IV.10.1.2}.
Every locally constructible subset being locally quasi-constructible, the maps $\mathfrak{C}(Y)\to\mathfrak{C}(X)$ and $\mathfrak{Lc}(Y)\to\mathfrak{Lc}(X)$ are injective; furthermore, for every open $U\subset Y$, the restriction $f^{-1}(U)\to U$ of $f$ is a quasi-homeomorphism \sref{IV.10.2.4}, hence if we show that $\mathfrak{C}(Y)\to\mathfrak{C}(X)$ is surjective, the same will be true of $\mathfrak{Lc}(Y)\to\mathfrak{Lc}(X)$.
But by virtue of \sref{IV.10.2.6}, every retrocompact open subset $Z$ in $X$ is of the form $f^{-1}(T)$, where $T$ is retrocompact open in $Y$; this evidently proves the surjectivity of $\mathfrak{C}(Y)\to\mathfrak{C}(X)$.
\end{proof}

\begin{remarks}[10.2.8]
\label{IV.10.2.8}
(i) We saw in the proof of \sref{IV.10.2.1} that if $f$ is a quasi-homeomorphism, the canonical map
\begin{equation}
\label{IV.10.2.8.1}
\tag{10.2.8.1}
\Gamma(Y,\sh{F})\to\Gamma(X,f^*(\sh{F}))
\end{equation}
is a functorial isomorphism in $\sh{F}$ in the category of sheaves of abelian groups on $Y$.
As $f^*$ is exact in this category \sref[0]{0.3.7.2}, this implies that the canonical homomorphisms $(T,3.2.2)$
\[
  \mathrm{H}^i(Y,\sh{F})\xrightarrow{\sim}\mathrm{H}^i(X,f^*(\sh{F}))
\]
are \emph{bijective} for every $i$.

(ii) If $f:X\to Y$ is a continuous map and $\sh{F}$, $\sh{G}$ two sheaves of abelian groups on $Y$, we have $f^*(\sh{F}\otimes_{\mathbf{Z}_Y}\sh{G})=f^*(\sh{F})\otimes_{\mathbf{Z}_X}f^*(\sh{G})$, an isomorphism canonical up to sign \sref[0]{0.4.3.3}.
We conclude that if $f$ is a quasi-homeomorphism, $f^*$ is also an \emph{equivalence} of the category of sheaves of rings on $Y$, and of the category of sheaves of rings on $X$; the data of a structure of ringed space on $Y$ is therefore equivalent to that of a structure of ringed space on $X$.

Given two ringed spaces $(X,\sh{O}_X)$, $(Y,\sh{O}_Y)$, we will say that a morphism of ringed spaces $f=(\psi,\theta):(X,\sh{O}_X)\to(Y,\sh{O}_Y)$ is a \emph{quasi-isomorphism} if $\psi$ is a quasi-homeomorphism of $X$ into $Y$ and $\theta^\sharp:\psi^*(\sh{O}_Y)\to\sh{O}_X$ an isomorphism of sheaves of rings; when this is so, the ringed space $(X,\sh{O}_X)$ is entirely determined, up to isomorphism, by $(Y,\sh{O}_Y)$, the space $X$, and the quasi-homeomorphism $\psi$ (which can be chosen arbitrarily).
When $f$ is a quasi-isomorphism of ringed spaces, the functor
\[
  \sh{F}\mapsto f^*(\sh{F})
\]
is an \emph{equivalence} of the category of $\sh{O}_Y$-Modules and of that of $\sh{O}_X$-Modules, since $f^*(\sh{F})$ is canonically identified here with $\psi^*(\sh{F})$.
We conclude for example isomorphisms of bi-$\partial$-functors
\[
  \operatorname{Ext}^i_{\sh{O}_Y}(\sh{F},\sh{G})\xrightarrow{\sim}\operatorname{Ext}^i_{\sh{O}_X}(f^*(\sh{F}),f^*(\sh{G})).
\]

In general, one can say that the usual constructions of the theory of sheaves and of homological algebra, carried out on the ringed space $Y$ or the ringed space $X$, are equivalent.
\end{remarks}

\subsection{Jacobson spaces}
\label{subsection:IV.10.3}

\oldpage[IV-3]{101}

\begin{definition}[10.3.1]
\label{IV.10.3.1}
We say that a topological space $X$ is a \emph{Jacobson space} if the set $X_0$ of closed points of $X$ is very dense in $X$ (in other words, if the canonical injection $X_0\to X$ is a quasi-homeomorphism).
\end{definition}

This means therefore \sref{IV.10.1.2} that every locally closed (or only locally quasi-constructible) subset $Z\neq\varnothing$ of $X$ contains a closed point of $X$, or equivalently that \emph{every closed subset of $X$ is the closure of the set of its closed points}.

\begin{proposition}[10.3.2]
\label{IV.10.3.2}
Let $X$ be a Jacobson space, $Z$ a quasi-constructible subset of $X$; then the subspace $Z$ is a Jacobson space, and for a point of $Z$ to be closed in $Z$, it is necessary and sufficient that it be closed in $X$.
\end{proposition}

\begin{proof}
If $X_0$ is the set of closed points of $X$, $Z\cap X_0$ is very dense in $Z$ by virtue of \sref{IV.10.2.4} applied to the injection $i:X_0\to X$; as the set $Z_0$ of closed points of $Z$ evidently contains $Z\cap X_0$, it is \emph{a fortiori} very dense in $Z$, hence $Z$ is a Jacobson space.
Let us now show that we have $Z_0=Z\cap X_0$; let $x\in Z$ be a closed point in $Z$; let $\overline{\{x\}}$ be its closure in $X$ and $\overline{\{x\}}\cap Z=\{x\}$, which is therefore a locally quasi-constructible subset of $X$, and as its intersection with $X_0$ is not empty \sref{IV.10.1.2}, we have $x\in X_0$.
\end{proof}

\begin{proposition}[10.3.3]
\label{IV.10.3.3}
Let $X$ be a topological space, $(U_\alpha)$ an open covering of $X$.
For $X$ to be a Jacobson space, it is necessary and sufficient that each of the subspaces $U_\alpha$ be a Jacobson space.
\end{proposition}

\begin{proof}
The condition is necessary by virtue of \sref{IV.10.3.2}.
Conversely, let us show first that the hypothesis that the $U_\alpha$ are Jacobson spaces implies that for a point $x\in U_\alpha$ to be closed in $X$, it suffices that it be closed in $U_\alpha$.
It suffices indeed to see that this condition implies that $x$ is also closed in each of the $U_\beta$ that contain it; but $U_\alpha\cap U_\beta$ is open in $U_\alpha$, hence $x$ is closed in $U_\alpha\cap U_\beta$, and by virtue of \sref{IV.10.3.2}, $x$ is also closed in $U_\beta$, which completes the proof.
\end{proof}

\subsection{Jacobson preschemes and Jacobson rings}
\label{subsection:IV.10.4}

\begin{definition}[10.4.1]
\label{IV.10.4.1}
We say that a prescheme $X$ is a \emph{Jacobson prescheme} if the underlying topological space of $X$ is a Jacobson space.
We say that a ring $A$ is a \emph{Jacobson ring} if $\operatorname{Spec}(A)$ is a Jacobson prescheme.
\end{definition}

Every closed subset of $X=\operatorname{Spec}(A)$ is of the form $Z=V(\mathfrak{a})$, where $\mathfrak{a}$ is an ideal equal to its radical, and the set $X_0$ of closed points of $X$ is the set of maximal ideals of $A$; to say that $Z\cap X_0$ is dense in $Z$ means therefore that $\mathfrak{a}$ is an \emph{intersection of maximal ideals} \sref[I]{I.1.1.4}; as $\mathfrak{a}$ is an intersection of prime ideals, this amounts to saying that every prime ideal of $A$ is an intersection of maximal ideals; by virtue of \sref{IV.10.3.1} and \sref{IV.10.1.2}, the usual definition of Jacobson rings (Bourbaki, \emph{Alg.\ comm.}, chap.~V, \S\,3, n$^\circ$\,4, def.~1) thus coincides with the definition \sref{IV.10.4.1}.

\begin{proposition}[10.4.2]
\label{IV.10.4.2}
\oldpage[IV-3]{102}
Let $X$ be a prescheme, $(U_\alpha)$ a covering of $X$ formed by affine open subsets.
For $X$ to be a Jacobson prescheme, it is necessary and sufficient that the rings $\Gamma(U_\alpha,\sh{O}_X)$ of the $U_\alpha$ be Jacobson rings.
\end{proposition}

This follows from \sref{IV.10.3.3} and the definition \sref{IV.10.4.1}.

\begin{env}[10.4.3]
\label{IV.10.4.3}
A discrete prescheme is a Jacobson prescheme; an Artinian ring is therefore a Jacobson ring.
Every principal ideal ring having an infinity of maximal ideals (for example $\mathbf{Z}$) is a Jacobson ring; a Noetherian local ring $A$ is a Jacobson ring if and only if its maximal ideal is its only prime ideal, i.e.\ if $A$ is Artinian.
Every sub-prescheme of a Jacobson prescheme is a Jacobson prescheme by virtue of \sref{IV.10.3.2}.
\end{env}

\begin{proposition}[10.4.4]
\label{IV.10.4.4}
Let $B$ be an integral ring.
The following conditions are equivalent:
\begin{enumerate}[label=\emph{\alph*)}]
  \item There exists $f\neq 0$ in $B$ such that $B_f$ is a field.
  \item The field of fractions of $B$ is a $B$-algebra of finite type.
  \item There exists a field $K$ containing $B$ and which is a $B$-algebra of finite type.
  \item The generic point of $\operatorname{Spec}(B)$ is isolated.
\end{enumerate}
\end{proposition}

\begin{proof}
It is clear that $d)$ is equivalent to $a)$, since $d)$ means that there exists $f\neq 0$ in $B$ such that $D(f)$ is reduced to the generic point of $\operatorname{Spec}(B)$.
It is trivial that $a)$ implies $b)$ and that $b)$ implies $c)$.
Finally, $c)$ implies $a)$, by virtue of Bourbaki, \emph{Alg.\ comm.}, chap.~V, \S\,3, n$^\circ$\,1, cor.~2 of th.~1.
\end{proof}

\begin{proposition}[10.4.5]
\label{IV.10.4.5}
Let $A$ be a ring.
The following conditions are equivalent:
\begin{enumerate}[label=\emph{\alph*)}]
  \item $A$ is a Jacobson ring.
  \item For every non-maximal prime ideal $\mathfrak{p}$ of $A$ and every $f\neq 0$ in $B=A/\mathfrak{p}$, $B_f$ is not a field.
  \item[\rm b$'$)] Every $A$-algebra of finite type $K$ which is a field is a finite $A$-algebra.
\end{enumerate}
\end{proposition}

\begin{proof}
We know that $a)$ implies $b')$ (Bourbaki, \emph{Alg.\ comm.}, chap.~V, \S\,3, n$^\circ$\,4, cor.~3 of th.~3).
Furthermore, the kernel of the homomorphism $A\to K$ is then a \emph{maximal} ideal $\mathfrak{m}$ of $A$, and $K$ is a finite extension of $A/\mathfrak{m}$ (\emph{loc.\ cit.}).
It is trivial that $b')$ implies $b)$, since $A/\mathfrak{p}=B$ is not a field and $B_f$ is a $B$-algebra of finite type.
It remains to show that $b)$ implies $a)$.
We will use the following lemma:
\end{proof}

\begin{lemma}[10.4.5.1]
\label{IV.10.4.5.1}
Let $X$ be a topological space having the following property: for every locally closed subset $Z\neq\varnothing$ of $X$, there exists a subset $Z'$ of $Z$, locally closed in $X$ (or in $Z$, which amounts to the same), and a point $x\in Z'$, closed in $Z'$.
Then, for $X$ to be a Jacobson space, it is necessary and sufficient that no non-closed point $x$ of $X$ be isolated in $\overline{\{x\}}$.
\end{lemma}

\begin{proof}
If $X$ is a Jacobson space, a non-closed point $x$ of $X$ cannot be isolated in $\overline{\{x\}}$, for this would mean that there exists an open $U$ of $X$ such that $U\cap\overline{\{x\}}=\{x\}$; but $U\cap\overline{\{x\}}$ is locally closed in $X$ and would contain no closed point in $X$, which contradicts the hypothesis that $X$ is a Jacobson space.
Conversely, suppose the conditions of the statement are satisfied; then the set $X_0$ of closed points of $X$ is identical to the set of $x\in X$ that are isolated in $\overline{\{x\}}$; but it follows from \sref{IV.5.1.10.1} that
\oldpage[IV-3]{103}
this set is \emph{very dense} in $X$, hence $X$ is a Jacobson space by definition.
\end{proof}

Let us recall \sref{IV.5.1.10} that the hypothesis made on $X$ in \sref{IV.10.4.5.1} is always satisfied when the underlying space of $X$ is a \emph{prescheme}.

Returning to the proof of \sref{IV.10.4.5}, condition $b)$ means that for every non-closed point $x$ of $\operatorname{Spec}(A)$, the generic point $x$ of $\operatorname{Spec}(A/\mathfrak{j}_x)=\overline{\{x\}}$ is not isolated in $\overline{\{x\}}$, hence $b)$ implies $a)$ by the lemma \sref{IV.10.4.5.1} and \sref{IV.10.4.4}.

\begin{corollary}[10.4.6]
\label{IV.10.4.6}
Every algebra of finite type $B$ over a Jacobson ring $A$ is a Jacobson ring, and the inverse image in $A$ of every maximal ideal of $B$ is a maximal ideal of $A$.
In particular, every algebra of finite type over a field or over $\mathbf{Z}$ is a Jacobson ring.
\end{corollary}

\begin{proof}
A $B$-algebra of finite type $K$ which is a field is also an $A$-algebra of finite type, hence an $A$-module of finite type and \emph{a fortiori} a $B$-module of finite type, whence the first assertion; the second was proved in the course of the proof of \sref{IV.10.4.5}, applied to $K=B/\mathfrak{m}$.
\end{proof}

\begin{corollary}[10.4.7]
\label{IV.10.4.7}
If $X$ is a Jacobson prescheme, $f:Y\to X$ a morphism locally of finite type, then $Y$ is a Jacobson prescheme and the image by $f$ of every closed point in $Y$ is a closed point in $X$.
\end{corollary}

\begin{proof}
The question being local on $X$ and on $Y$, one reduces to the case where $X$ and $Y$ are affine, and the corollary then follows from \sref{IV.10.4.6}.
\end{proof}

\begin{corollary}[10.4.8]
\label{IV.10.4.8}
If $k$ is an algebraically closed field, $X$ a $k$-prescheme locally of finite type, the set of rational points of $X$ over $k$ is very dense in $X$.
\end{corollary}

\begin{proof}
Indeed, $X$ is a Jacobson prescheme \sref{IV.10.4.7} and the closed points are exactly the rational points over $k$ \sref[I]{I.6.4.2}.
\end{proof}

\begin{env}[10.4.9]
\label{IV.10.4.9}
The fact that preschemes locally of finite type over a field or over $\mathbf{Z}$ are Jacobson preschemes is particularly important, because of the possibility of reducing to this case in many questions of algebraic geometry \sref{IV.8.1.2},~$c)$.
We will give two examples:
\end{env}

\begin{env}[10.4.10]
\label{IV.10.4.10}
\textbf{Applications (10.4.10) :~I~: Proof of \sref{IV.6.15.9}.}
Let $k$ be a separably closed field, $X$ a $k$-prescheme locally of finite type over $k$ and \emph{unibranch}.
We know that the integral closure of a $k$-algebra of finite type $A$ in an extension of finite type of its field of fractions is a finite $A$-algebra (Bourbaki, \emph{Alg.\ comm.}, chap.~V, \S\,3, n$^\circ$\,2, th.~2), hence every $k$-algebra of finite type is a universally Japanese ring; we conclude that the set of points $x\in X$ where $X$ is geometrically unibranch is locally constructible \sref{IV.9.7.10}.
But the hypothesis and the lemma \sref{IV.6.15.8} imply that this set contains all the \emph{closed} points of $X$.
The conclusion therefore follows from \sref{IV.10.4.6}, \sref{IV.10.3.1} and the bijectivity of the canonical map $\mathfrak{Lc}(X)\to\mathfrak{Lc}(X_0)$ (where $X_0$ is the set of closed points of $X$) \sref{IV.10.2.7}.
\end{env}

\begin{env}[10.4.11]
\label{IV.10.4.11}
\oldpage[IV-3]{104}
\textbf{Applications :~II~: Proposition (10.4.11).}
Let $S$ be a prescheme, $X$ an $S$-prescheme of finite type.
Every $S$-endomorphism of $X$ which is radicial is surjective (hence bijective).
\end{env}

\begin{proof}
Let $f:X\to S$ be the structural morphism, $g:X\to X$ the $S$-endomorphism considered, and, for every $s\in S$, let $g_s$ be the morphism deduced from $g$ by the base change $\operatorname{Spec}(k(s))\to S$, which is a $k(s)$-endomorphism of the fibre $f^{-1}(s)=X\times_S\operatorname{Spec}(k(s))$.

To prove that $g$ is surjective, it suffices to prove that $g_s$ is surjective for every $s\in S$, hence one can (by virtue of \sref[I]{I.3.5.7}) reduce to the case where $S=\operatorname{Spec}(k)$ is the spectrum of a field $k$, in which case $f$ is a morphism of finite presentation, since $S$ is Noetherian.
Applying \sref{IV.8.9.1} and \sref{IV.8.10.5},~(vii), one reduces to the case where $S=\operatorname{Spec}(A)$, where $A$ is a sub-$\mathbf{Z}$-algebra of finite type of $k$.
But then $X$ is a Jacobson prescheme \sref{IV.10.4.7}, and $g(X)$ is constructible in $X$ \sref[I]{I.1.8.5}, hence, to prove that $g(X)=X$, it suffices to show that $g(X)$ contains all the \emph{closed} points of $X$ \sref{IV.10.3.1}.

\begin{lemma}[10.4.11.1]
\label{IV.10.4.11.1}
Let $Y$ be a $\mathbf{Z}$-prescheme of finite type.
\begin{enumerate}[label=(\roman*)]
  \item For a closed point $y\in Y$, it is necessary and sufficient that $k(y)$ be a finite field.
  \item For every prime number $p$ and every integer $d\geq 1$, the set of points $y\in Y$ such that $k(y)$ is an extension of $\mathbf{F}_p$ of degree dividing $d$ is finite.
\end{enumerate}
\end{lemma}

The assertion (i) follows from the fact that the image of a closed point $y\in Y$ in $\operatorname{Spec}(\mathbf{Z})$ is a closed point \sref{IV.10.4.7}, in other words a prime number, and from \sref[I]{I.6.4.2}.
On the other hand, as $Y$ is a finite union of affine open subsets of finite type over $\mathbf{Z}$, one can restrict to proving (ii) in the case where $Y=\operatorname{Spec}(C)$, where $C$ is a $\mathbf{Z}$-algebra of finite type.
Now, the points $y\in Y$ such that the degree of $k(y)$ over $\mathbf{F}_p$ divides $d$ correspond bijectively to the homomorphisms $C\to\mathbf{F}_{p^d}$; but if $(t_i)$ is a system of generators of the $\mathbf{Z}$-algebra $C$, every homomorphism of $C$ is determined by its values on the elements $t_i$, and there are therefore only a finite number of homomorphisms of $C$ into a finite field.

This being so, since $X$ is a $\mathbf{Z}$-prescheme of finite type, let $T_{p,d}$ be the set of closed points $z\in X$ such that $k(z)$ is an extension of $\mathbf{F}_p$ of degree dividing $d$; it follows from \sref{IV.10.4.11.1} that the set $T_{p,d}$ is finite and that the set of closed points of $X$ is the union of the $T_{p,d}$.
Furthermore, if $z\in T_{p,d}$ and if $h$ is any endomorphism of $X$, $k(h(z))$ is isomorphic to a subfield of $k(z)$, hence $h(z)\in T_{p,d}$; in other words $T_{p,d}$ is \emph{stable} under every endomorphism of $X$.
As $g$ is by hypothesis injective, its restriction to $T_{p,d}$ is \emph{bijective} since $T_{p,d}$ is finite, which completes proving the proposition.

We will see further on \sref{IV.17.9.6} that when one further supposes that $X$ is an $S$-prescheme of finite presentation, and on the other hand that $g$ is a monomorphism, then one can assert that $g$ is an \emph{automorphism} of $X$.
\end{proof}

\subsection{Noetherian Jacobson preschemes}
\label{subsection:IV.10.5}

\begin{proposition}[10.5.1]
\label{IV.10.5.1}
Let $B$ be a Noetherian integral ring.
The conditions equivalent to $a)$ through $d)$ of \sref{IV.10.4.4} are then also equivalent to the following:
\begin{enumerate}[label=\emph{\alph*)}]
  \setcounter{enumi}{4}
  \item $\operatorname{Spec}(B)$ is finite.
  \item $B$ is a semi-local ring of dimension $\leq 1$.
\end{enumerate}
\end{proposition}

\begin{proof}
\oldpage[IV-3]{105}
It follows indeed from the Artin--Tate theorem \sref[0]{0.16.3.3} that conditions $e)$ and $f)$ are equivalent.
Condition $f)$ implies that $\operatorname{Spec}(B)$ is the union of the set of its closed points and of its generic point, hence $f)$ implies $e)$ without assuming $A$ to be Noetherian.
Finally, $e)$ implies $d)$ without assuming $A$ to be Noetherian, for the generic point $x$ of $X$ is the complement of the union of the closures $\overline{\{y\}}$, where $y$ ranges over the set of points $y\neq x$, and as these points are finite in number, $\{x\}$ is the complement of a closed set in $X$.
\end{proof}

\begin{corollary}[10.5.2]
\label{IV.10.5.2}
For a Noetherian ring $A$ to be a Jacobson ring, it is necessary and sufficient that there exist no prime ideal $\mathfrak{p}$ of $A$ such that $A/\mathfrak{p}$ is a semi-local ring of dimension $1$.
\end{corollary}

\begin{proof}
This follows immediately from \sref{IV.10.5.1} and from condition $b)$ of \sref{IV.10.4.5}, the prime ideals $\mathfrak{p}$ of $A$ such that $A/\mathfrak{p}$ is semi-local of dimension~$0$ being the maximal ideals of $A$.
\end{proof}

\begin{corollary}[10.5.3]
\label{IV.10.5.3}
Let $X$ be an irreducible Noetherian prescheme.
The following conditions are equivalent:
\begin{enumerate}[label=\emph{\alph*)}]
  \item The generic point of $X$ is isolated.
  \item $X$ is finite.
  \item $X$ is of dimension $\leq 1$ and the set of its closed points is finite.
\end{enumerate}
\end{corollary}

\begin{proof}
There exists by hypothesis a finite number of irreducible affine open subsets $U_i$ ($1\leq i\leq n$) covering $X$, and each of which contains the generic point of $X$; it suffices to prove the equivalence of $a)$, $b)$, and $c)$ for each of the $(U_i)_{\mathrm{red}}$ (taking into account \sref[0]{0.14.1.7}).
But this equivalence follows from \sref{IV.10.5.1}.
\end{proof}

\begin{remark}[10.5.4]
\label{IV.10.5.4}
A Noetherian prescheme $X$ satisfying the equivalent conditions of \sref{IV.10.5.3} is not necessarily an affine scheme; in fact it can even be \emph{non-separated}, as one sees by replacing, in the example \sref[I]{I.5.5.11} of the ``affine line with doubled point'', $X_1$ and $X_2$ by the spectrum of the ring of the discrete valuation $(K[s])_{(s)}$, $U_{12}$ and $U_{21}$ by the open subset reduced to the generic point in $X_1$ and $X_2$ respectively; the non-separated prescheme $X$ that one obtains has exactly $3$ points.
\end{remark}

\begin{proposition}[10.5.5]
\label{IV.10.5.5}
Let $X$ be a Noetherian prescheme.
\begin{enumerate}[label=(\roman*)]
  \item The set $X_0$ of points $x\in X$ such that $\overline{\{x\}}$ is finite is very dense in $X$.
  \item For $X$ to be a Jacobson prescheme, it is necessary and sufficient that there exist no sub-prescheme of $X$ isomorphic to the spectrum of a semi-local integral ring of dimension $1$.
\end{enumerate}
\end{proposition}

\begin{proof}
The condition that $\overline{\{x\}}$ is finite is equivalent, here \sref{IV.10.5.3}, to the fact that $x$ is isolated in $\overline{\{x\}}$, it being the underlying space of a sub-prescheme (Noetherian) of $X$; the assertion (i) therefore follows from \sref{IV.10.4.5.1}.
Similarly, taking into account \sref{IV.10.4.5.1}, to show the assertion (ii), let us note that for a non-closed point $x$ of $X$ to belong to $X_0$, it is necessary and sufficient that the closed sub-prescheme $Y$ of $X$ having $\overline{\{x\}}$ as underlying space be of dimension $1$ and finite, hence a finite union of sub-preschemes which are open (in $Y$) and affine $U_i$ which are spectra of semi-local integral rings of dimension $1$.
Conversely, if there is a sub-prescheme $Z$ of $X$ which is the spectrum of a semi-local integral ring of dimension $1$, $Z$ is not a Jacobson prescheme \sref{IV.10.5.2}, hence neither is $X$ \sref{IV.10.3.2}.
\end{proof}

\oldpage[IV-3]{106}

\begin{remark}[10.5.6]
\label{IV.10.5.6}
The assertion (ii) of \sref{IV.10.5.5} is still valid when $X$ is \emph{locally Noetherian}: indeed, if $(V_\alpha)$ is a covering of $X$ by affine open subsets (Noetherian), every sub-prescheme of one of the $V_\alpha$ is a sub-prescheme of $X$; conversely, if a sub-prescheme $Z$ of $X$ is isomorphic to the spectrum of a semi-local integral ring of dimension $1$, there is an $\alpha$ such that $V_\alpha\cap Z$ contains an affine open $U$ of $Z$ which is non-reduced to the generic point of $Z$, and which is therefore also the spectrum of a semi-local integral ring of dimension $1$.
We conclude by \sref{IV.10.3.3}.
\end{remark}

\begin{proposition}[10.5.7]
\label{IV.10.5.7}
Let $X$ be a locally Noetherian prescheme, $Y$ a closed subset of $X$ such that every non-empty closed subset of $X$ meets $Y$.
Then the subprescheme induced on the open subset $X-Y$ is a Jacobson prescheme.
\end{proposition}

\begin{proof}
Applying the criterion \sref{IV.10.5.5},~(ii), and supposing that there were in $X-Y$ a sub-prescheme $Z$ which is the spectrum of a semi-local integral ring of dimension $1$, the generic point $z$ of $Z$ being isolated in $Z$ (or in the closure $\overline{Z}$ of $Z$ in $X$), and $Z$ being distinct from $\{z\}$.
Let $y\neq z$ be a point of $Z$; as it does not belong to $Y$, it is not a closed point in $X$, and its closure $\overline{\{y\}}$ in $X$ meets $Y$ at a point $x\neq y$, which is therefore a specialisation of $y$.
The existence of the chain $\{x\}\subset\overline{\{y\}}\subset\overline{\{z\}}$ then shows that the dimension of $\overline{\{z\}}$ would be $\geq 2$, and the same would be true of the dimension of $U\cap\overline{\{z\}}$, where $U$ is an affine (hence Noetherian) open neighbourhood of $x$ in $X$.
But this contradicts the fact that the generic point $z$ of $U\cap\overline{\{z\}}$ is isolated in $U\cap\overline{\{z\}}$ \sref{IV.10.5.3}.
\end{proof}

\begin{corollary}[10.5.8]
\label{IV.10.5.8}
Let $A$ be a Noetherian ring; for every element $f$ of the radical $\mathfrak{R}$ of $A$, the ring $A_f$ is a Jacobson ring, and the open subset $\operatorname{Spec}(A)-V(\mathfrak{R})$ is a Jacobson scheme.
\end{corollary}

\begin{proof}
If $X=\operatorname{Spec}(A)$, $Y=V(\mathfrak{J})$, where $\mathfrak{J}$ is an ideal of $A$, to say that every non-empty closed subset of $X$ meets $Y$ is equivalent \sref[0]{0.2.1.3} to saying that $Y$ contains all the \emph{closed points} of $X$, or equivalently that $\mathfrak{J}$ is \emph{contained in the radical} $\mathfrak{R}$ of $A$.
If $f\in\mathfrak{R}$, the open $D(f)$ does not meet $V(\mathfrak{R})$, and is a Jacobson space by virtue of \sref{IV.10.5.7}.
\end{proof}

\begin{corollary}[10.5.9]
\label{IV.10.5.9}
Let $A$ be a Noetherian local ring, $\mathfrak{m}$ its maximal ideal, $X=\operatorname{Spec}(A)$, $Y=X-\{\mathfrak{m}\}$; then $Y$ is a Jacobson scheme, whose closed points are the prime ideals $\mathfrak{p}\in\operatorname{Spec}(A)$ such that $\dim(A/\mathfrak{p})=1$.
\end{corollary}

\begin{proof}
The first assertion is a particular case of \sref{IV.10.5.8}; on the other hand the closed points of $Y$ are the prime ideals $\mathfrak{p}$ of $A$ which are maximal elements in the set of prime ideals $\neq\mathfrak{m}$, which, by definition of the dimension, means that $\dim(A/\mathfrak{p})=1$.
\end{proof}

\begin{proposition}[10.5.10]
\label{IV.10.5.10}
Let $A$ be a Noetherian local ring, reduced, complete, and not a field.
Then the finite intersections of the kernels of the local homomorphisms of $A$ into discrete valuation rings $V$, which are $A$-algebras of finite type, form a basis of filter tending to $0$ for the adic topology of $A$.
\end{proposition}

\begin{proof}
\oldpage[IV-3]{107}
It suffices (Bourbaki, \emph{Alg.\ comm.}, chap.~III, \S\,2, n$^\circ$\,7, prop.~8) to prove that the intersection of the kernels considered in the statement is reduced to $0$.
Suppose first that $A$ is integral and of dimension $1$; by virtue of the Nagata theorem \sref[0]{0.23.1.5} and \sref{0.23.1.6}, the integral closure $A'$ of $A$ is a complete local integral ring and
integrally closed and of dimension $1$; it is therefore a discrete valuation ring, and the proposition follows immediately in this case.

Let us pass to the general case; let $x$ be the closed point of $X=\operatorname{Spec}(A)$, and set $Y=X-\{x\}$; we know \sref{IV.10.5.9} that $Y$ is a Jacobson prescheme.
Let us show that the intersection of the prime ideals $\mathfrak{p}$ of $A$ such that $\dim(A/\mathfrak{p})=1$ is reduced to $0$.
But to say that the intersection of these prime ideals is reduced to $0$ means that the set of these ideals is dense in $X$, or equivalently in $Y$ (since $Y$ is dense in $X$), and this follows immediately from \sref{IV.10.5.9}.
\end{proof}

\subsection{Dimension in Jacobson preschemes}
\label{subsection:IV.10.6}

The results of this section give precision in certain cases and generalise results from \S\,5.

\begin{proposition}[10.6.1]
\label{IV.10.6.1}
Let $S$ be a locally Noetherian prescheme, satisfying in addition the conditions: $1^\circ$~$S$ is a Jacobson prescheme; $2^\circ$~for every $s\in S$, $\sh{O}_{S,s}$ is universally catenary \sref{IV.5.6.2}; $3^\circ$~every irreducible component $S'$ of $S$ is equicodimensional (in other words, for every closed point $s$ of $S'$ and every sub-prescheme $S'$ having $S'$ as underlying space, we have $\dim(S')=\dim(\sh{O}_{S',s})$).
We then have the following properties:
\begin{enumerate}[label=(\roman*)]
  \item For every morphism locally of finite type $g:X\to S$, $X$ satisfies conditions $1^\circ$, $2^\circ$, and $3^\circ$ above.
In particular, if $X$ is equidimensional (for example if $X$ is irreducible), $X$ is biequidimensional (in other words $X$ is catenary and for every closed point $x$ of $X$, we have $\dim(\sh{O}_{X,x})=\dim(X)$ \sref[0]{0.14.3.3}).
  \item Let $X$, $Y$ be two $S$-preschemes locally of finite type over $S$, $f:X\to Y$ an $S$-morphism; suppose $X$ irreducible and $f$ dominant.
If $\xi$ (resp.\ $\eta$) is the generic point of $X$ (resp.\ $Y$) and $e=\dim(f^{-1}(\eta))=\deg.\mathrm{tr}_{k(\eta)}k(\xi)$, we have
\begin{equation}
\label{IV.10.6.1.1}
\tag{10.6.1.1}
\dim(X)=\dim(Y)+e.
\end{equation}
  \item Let $X$, $Y$ be two $S$-preschemes locally of finite type over $S$, $f:X\to Y$ an $S$-morphism, $n$ a non-negative integer such that $\dim(f^{-1}(y))\geq n$ (resp.\ $\dim(f^{-1}(y))\leq n$) for every $y\in Y$.
Then we have
\begin{equation}
\label{IV.10.6.1.2}
\tag{10.6.1.2}
\dim(X)\geq\dim(Y)+n
\end{equation}
(resp.
\begin{equation}
\label{IV.10.6.1.3}
\tag{10.6.1.3}
\dim(X)\leq\dim(Y)+n).
\end{equation}
\end{enumerate}
\end{proposition}

\begin{proof}
\oldpage[IV-3]{108}
(i) The property $1^\circ$ for $X$ follows from \sref{IV.10.4.7}.
For every $x\in X$, $\sh{O}_{X,x}$ is the local ring of an $\sh{O}_{S,g(x)}$-algebra of finite type at a prime ideal, and the homomorphism $\sh{O}_{S,g(x)}\to\sh{O}_{X,x}$ is local; hence \sref{IV.5.6.3},~(iv) $\sh{O}_{X,x}$ is universally catenary.
To show that $X$ satisfies condition $3^\circ$, consider several cases:

\emph{a)} $X$ is a closed irreducible sub-prescheme of $S$; let $S'$ be an irreducible component of $S$ containing $X$, $\xi$ the generic point of $X$, $x$ a closed point of $X$; for every $s\in S'$, $\sh{O}_{S',s}$, a quotient of $\sh{O}_{S,s}$, is catenary \sref{IV.5.6.1}, hence conditions $2^\circ$ and $3^\circ$ imply that $S'$ is biequidimensional \sref[0]{0.14.3.3}.
By virtue of \sref{IV.5.1.2} and of \sref[0]{0.14.3.3.2}, we therefore have
\[
  \dim(\sh{O}_{X,x})=\dim(\sh{O}_{S',x})-\dim(\sh{O}_{S',\xi})=\dim(S')-\dim(\sh{O}_{S',\xi})
\]
which proves that $\dim(\sh{O}_{X,x})$ does not depend on the closed point $x$ considered, whence the assertion in this case \sref{IV.5.1.4}.

\emph{b)} $X$ is irreducible and $g$ dominant.
Then, for every closed point $x$ of $X$, $g(x)$ is closed in $S$ by \sref{IV.10.4.7}; as $\sh{O}_{S,g(x)}$ is universally catenary, it follows from \sref{IV.5.6.5.3} that we have
\[
  \dim(\sh{O}_{X,x})=\dim(\sh{O}_{S,g(x)})+e
\]
where $e=\dim(g^{-1}(\zeta))$, $\zeta$ being the generic point of $S$.
As $\dim(S)=\dim(\sh{O}_{S,g(x)})$ by virtue of condition $3^\circ$ for $S$, we have $\dim(\sh{O}_{X,x})=\dim(S)+e$ for every closed point $x\in X$; this proves condition $3^\circ$ for $X$ \sref{IV.5.1.4}, and at the same time the formula
\begin{equation}
\label{IV.10.6.1.4}
\tag{10.6.1.4}
\dim(X)=\dim(S)+e.
\end{equation}

\emph{c) General case.}
Considering a reduced sub-prescheme $X'$ of $X$ having for underlying space an irreducible component of $X$, one reduces to the case where $X$ is integral; using \sref[I]{I.5.2.2} and the case $a)$ proved above, one can then replace $S$ by the reduced sub-prescheme having $g(X)$ as underlying space; one is then brought to case $b)$, and this completes proving (i).

(ii) The morphism $f$ being locally of finite type \sref[I]{I.1.3.4}, one can apply the results of (i) replacing $S$ by $Y$; furthermore, as $X$ is irreducible and $f$ dominant, one can also replace $S$ by $Y$ in \sref{IV.10.6.1.4}, which gives \sref{IV.10.6.1.1}.

(iii) The relative assertion where $\dim(f^{-1}(y))\leq n$ for every $y\in Y$ has already been proved under more general hypotheses in \sref{IV.5.6.7}.
Suppose that $\dim(f^{-1}(y))\geq n$ for every $y\in Y$, and consider a generic point $\eta$ of an irreducible component $Y'$ of $Y$; there exists at least one irreducible component $Z$ of $f^{-1}(\eta)$ of dimension $\geq n$; if $\xi$ is the generic point of $Z$, $\xi$ is also a generic point of an irreducible component $X'$ of $X$ such that $f(X')$ is dense in $Y'$ \sref[0]{0.2.1.8}.
Consider the reduced sub-preschemes of $X$, $Y$ having respectively for underlying spaces $X'$, $Y'$, and the restriction $X'\to Y'$ of $f$ \sref[I]{I.5.2.2}; it then follows from (ii) that $\dim(X')\geq\dim(Y')+n$, and \emph{a fortiori} $\dim(X)\geq\dim(Y')+n$; this being true for every irreducible component $Y'$ of $Y$, we conclude the inequality \sref{IV.10.6.1.2}.
\end{proof}

\begin{corollary}[10.6.2]
\label{IV.10.6.2}
Suppose that $X$ satisfies conditions $1^\circ$, $2^\circ$, and $3^\circ$ of \sref{IV.10.6.1}.
Then, for every dense open $U$ in $X$, we have $\dim(U)=\dim(X)$.
\end{corollary}

\begin{proof}
\oldpage[IV-3]{109}
We know indeed that $\dim(U)\leq\dim(X)$ \sref[0]{0.14.1.4}; furthermore, as $X$ is a Jacobson prescheme, $U$ contains a closed point of every irreducible component of $X$, hence $\dim(U)=\dim(X)$ by virtue of \sref{IV.10.6.1} and \sref[0]{0.14.1.2.1}.
\end{proof}

\begin{proposition}[10.6.3]
\label{IV.10.6.3}
Suppose that $X$ satisfies conditions $1^\circ$, $2^\circ$, and $3^\circ$ of \sref{IV.10.6.1}, and let $Y$ be a closed subset of $X$.
For every $x\in Y$, and every open neighbourhood $U$ of $x$ in $X$ not meeting the irreducible components of $Y$ that do not contain $x$, we have
\begin{equation}
\label{IV.10.6.3.1}
\tag{10.6.3.1}
\dim_x(Y)=\dim(U\cap Y)=\dim\overline{\{x\}}+\operatorname{codim}(\overline{\{x\}},Y)=\sup_i\dim(Y_i)
\end{equation}
where $Y_i$ ($1\leq i\leq m$) are the irreducible components of $Y$ containing $x$.
\end{proposition}

\begin{proof}
Considering a closed sub-prescheme of $X$ having $Y$ as underlying space, one can restrict, by virtue of \sref{IV.10.6.1},~(i), to the case where $Y=X$.
By the choice of $U$, $U$ is the union of the $U\cap X_i$, hence $\dim(U)=\sup_i\dim(U\cap X_i)$; but by \sref{IV.10.6.2} applied to a closed sub-prescheme of $X$ having $X_i$ as underlying space, we have (taking into account \sref{IV.10.6.1},~(i)), $\dim(U\cap X_i)=\dim(X_i)$; this proves that the second and the fourth terms of \sref{IV.10.6.3.1} are equal.
As $U\cap X_i$ is biequidimensional \sref{IV.10.6.1},~(i) (since the immersion $U\cap X_i\to X_i$ is of finite type \sref[I]{I.6.3.5}), we have

\begin{equation}
\label{IV.10.6.3.2}
\tag{10.6.3.2}
\dim(X_i)=\dim(U\cap X_i)=\dim(U\cap\overline{\{x\}})+\operatorname{codim}(U\cap\overline{\{x\}},U\cap X_i)
\end{equation}
\sref[0]{0.14.3.5}.
By \sref{IV.10.6.2} applied to a closed sub-prescheme of $X$ having $\overline{\{x\}}$ as underlying space, $\dim(U\cap\overline{\{x\}})=\dim(\overline{\{x\}})$; as we also have, for the same reasons,

\begin{equation}
\label{IV.10.6.3.3}
\tag{10.6.3.3}
\dim(X_i)=\dim(\overline{\{x\}})+\operatorname{codim}(\overline{\{x\}},X_i)
\end{equation}
we obtain
\[
  \dim(U)=\dim(\overline{\{x\}})+\sup_i\operatorname{codim}(\overline{\{x\}},X_i)=\dim(\overline{\{x\}})+\operatorname{codim}(\overline{\{x\}},X)
\]
by definition of the codimension \sref[0]{0.14.2.1}.
This shows that $\dim(U)$ is \emph{independent} of the open neighbourhood $U$ of $x$ satisfying the conditions of the statement, hence is equal to $\dim_x(X)$, by \sref[0]{0.14.1.4.1}.
\end{proof}

\begin{corollary}[10.6.4]
\label{IV.10.6.4}
Under the hypotheses of \sref{IV.10.6.3}, let $\sh{F}$ be a coherent $\sh{O}_X$-Module, $Y$ its support.
For every $x\in Y$, we have
\begin{equation}
\label{IV.10.6.4.1}
\tag{10.6.4.1}
\dim(\overline{\{x\}})+\dim(\sh{F}_x)=\dim_x Y.
\end{equation}
\end{corollary}

\begin{proof}
Indeed, this follows from \sref{IV.10.6.3.1} and from the formula $\dim(\sh{F}_x)=\operatorname{codim}(\overline{\{x\}},Y)$ \sref{IV.5.1.12.2}.
\end{proof}

\subsection{Examples and counter-examples}
\label{subsection:IV.10.7}

\begin{env}[10.7.1]
\label{IV.10.7.1}
Let $S$ be a locally Noetherian prescheme of dimension $\leq 1$ and suppose that $S$ is a \emph{Jacobson prescheme}; when $S$ is Noetherian, this amounts to saying that the irreducible components of $S$ of dimension $1$ are infinite, for every $x\in S$ which is not a closed point is the generic point of such a component \sref{IV.10.4.5} and \sref{IV.10.5.4}.
Then $S$ also satisfies conditions $2^\circ$ and $3^\circ$ of \sref{IV.10.6.1}: indeed, every local ring $\sh{O}_{S,s}$ is of dimension $0$ or $1$, and hence universally catenary \sref{IV.5.6.3}, \sref{IV.7.2.9}; on the other hand, an irreducible component $S'$ of $S$ is, either reduced to a closed point $s$ of $S'$, or biequidimensional.
\end{env}

\oldpage[IV-3]{110}

From these remarks and from \sref{IV.10.6.1} it follows that every prescheme locally of finite type over $S$ also satisfies the properties $1^\circ$, $2^\circ$, and $3^\circ$; it is in particular the case of preschemes \emph{locally of finite type over a field or over $\mathbf{Z}$}.

\begin{env}[10.7.2]
\label{IV.10.7.2}
Let $A$ be a Noetherian local ring \emph{universally catenary} and let $S$ be the complementary of \emph{the closed point} $a$ in $X=\operatorname{Spec}(A)$.
Then $S$ satisfies conditions $1^\circ$, $2^\circ$, and $3^\circ$ of \sref{IV.10.6.1}: indeed, one has already seen that $S$ is a Jacobson prescheme \sref{IV.10.5.9}; as $A$ is universally catenary, so are all the local rings $A_\mathfrak{p}$ at the prime ideals of $A$ \sref{IV.5.6.3}.
On the other hand, an irreducible component $S'$ of $S$ is the complement of $a$ in an irreducible component $X'$ of $X$. For every closed point $x$ of $S$, the closure of $x$ in $X$ is $\{x,a\}$, so $\dim(\overline{\{x\}})=1$. Since $X'$ is biequidimensional \sref[0]{0.14.3.3}, it follows that $\dim(\sh{O}_{X',x})=\dim(X')-1$ \sref[0]{0.14.3.2}, proving that $S'$ is equicodimensional.
\end{env}

\begin{env}[10.7.3]
\label{IV.10.7.3}
Let $A$ be a discrete valuation ring; let us show that in $B=A[T_1,\ldots,T_n]$, with $n$ indeterminates, there exist two maximal ideals $\mathfrak{m}$, $\mathfrak{n}$ of \emph{respective heights} $n$ and $n+1$.
We have seen \sref{IV.5.2.5},~(i) for $n=1$; let us prove it by induction on $n$.
As $A[T_1,\ldots,T_{n+1}]$ is a free, hence faithfully flat, $A[T_1,\ldots,T_n]$-module, there are maximal ideals $\mathfrak{m}'$ and $\mathfrak{n}'$ of $A[T_1,\ldots,T_{n+1}]$ lying respectively over $\mathfrak{m}$ and $\mathfrak{n}$ \sref[0]{0.6.5.1}; by \sref{IV.5.5.3}, these ideals have respective heights $n+1$ and $n+2$, proving the assertion.
Suppose $n\geq 2$ in what follows.
Let $\mathfrak{J}=\mathfrak{m}\cap\mathfrak{n}=\mathfrak{m}\mathfrak{n}$ and let $R=1+\mathfrak{J}$, which is a multiplicative subset of $B$.
If $B'=R^{-1}B$, the ideal $\mathfrak{J}'=R^{-1}\mathfrak{J}$ is contained in the radical of $B'$ (Bourbaki, \emph{Alg\`ebre commutative}, Chap.~III, \S3, no.~5, Prop.~12).
As a topological space, $X'=\operatorname{Spec}(B')$ is identified with a subspace of $X=\operatorname{Spec}(B)$, and at the points $x$ of $X'$ the local rings $\sh{O}_{X,x}$ and $\sh{O}_{X',x}$ are the same \sref[I]{I.1.6.2}.
Consider in $X'$ the closed subset $Y'=V(\mathfrak{J}')$, and put $S=X'\setminus Y'$.
Then $S$ is a Jacobson prescheme \sref{IV.10.5.7}, and is evidently irreducible and Noetherian; moreover, for every $s\in S$, the local ring $\sh{O}_{S,s}=\sh{O}_{X,s}$ is universally catenary by \sref{IV.5.6.3}, since $A$ is universally catenary \sref{IV.5.6.4}.
Nevertheless, there are two closed points $a,b$ of $S$ such that $\sh{O}_{S,a}$ and $\sh{O}_{S,b}$ \emph{do not have the same dimension}; in other words, $S$ does not satisfy condition $3^\circ$ of \sref{IV.10.6.1}.
To see this, consider the two maximal ideals $\mathfrak{m}'=R^{-1}\mathfrak{m}$ and $\mathfrak{n}'=R^{-1}\mathfrak{n}$ of $B'$, which have heights $n$ and $n+1$, respectively.
We have $\mathfrak{J}'=\mathfrak{m}'\cap\mathfrak{n}'$, and hence $\mathfrak{J}'$ is contained in no prime ideal of $B'$ distinct from $\mathfrak{m}'$ and $\mathfrak{n}'$; these are consequently the only maximal ideals of $B'$.
Let $a',b'$ be the only closed points of $X'$, corresponding to $\mathfrak{m}'$ and $\mathfrak{n}'$.
There exists in $B'$ a nonmaximal prime ideal $\mathfrak{p}'\subset\mathfrak{m}'$ which is not contained in $\mathfrak{n}'$: it suffices to show that in $B$ there is a nonmaximal prime ideal contained in $\mathfrak{m}$ and not in $\mathfrak{n}$; for this one may, for example, consider the fibre of $\mathfrak{m}$ for the morphism corresponding to the injection $A[T_1]\longrightarrow B=A[T_1,\ldots,T_n]$ and apply \sref{IV.6.1.2}.
By considering a maximal chain of prime ideals between $\mathfrak{p}'$ and $\mathfrak{m}'$, and replacing $\mathfrak{p}'$ by the penultimate ideal in this chain, we may therefore suppose that the point $a$ of $X'$ corresponding to $\mathfrak{p}'$ has closure $\{a,a'\}$ in $X'$.
Since $B'_{\mathfrak{m}'}$ is biequidimensional, $\mathfrak{p}'$ then has height $n-1$.
In the same way one constructs a nonmaximal prime ideal $\mathfrak{q}'$ of $B'$ of height $n$ such that, if $b$ is the corresponding point of $X'$, the closure of $b$ in $X'$ is $\{b,b'\}$.
Thus $a$ and $b$ belong to $S$ and are closed in $S$, and consequently answer the question.
\end{env}

\subsection{Rectified depth}
\label{subsection:IV.10.8}

\begin{definition}[10.8.1]
\label{IV.10.8.1}
Let $X$ be a locally Noetherian prescheme, $\sh{F}$ a coherent $\sh{O}_X$-Module.
For every $x\in X$, we call \emph{rectified depth} of $\sh{F}$ at the point $x$ and denote $\operatorname{prof}^*_x(\sh{F})$ the number (integer $\geq 0$ or $+\infty$) equal to
\begin{equation}
\label{IV.10.8.1.1}
\tag{10.8.1.1}
\operatorname{prof}^*_x(\sh{F})=\operatorname{prof}(\sh{F}_x)+\dim(\overline{\{x\}})
\end{equation}
where $\overline{\{x\}}$ is the closure of the point $x$ in $X$.
For every closed subset $Z$ of $X$, we call \emph{rectified depth of $\sh{F}$ along $Z$} and denote $\operatorname{prof}^*_Z(\sh{F})$ the number
\begin{equation}
\label{IV.10.8.1.2}
\tag{10.8.1.2}
\operatorname{prof}^*_Z(\sh{F})=\inf_{x\in Z}\operatorname{prof}^*_x(\sh{F}).
\end{equation}
\end{definition}

In other terms, for every integer $n$, the relation $\operatorname{prof}^*_Z(\sh{F})\geq n$ is equivalent to $\operatorname{prof}^*_x(\sh{F})\geq n$ for every $x\in Z$.
If $Z=X$, one writes $\operatorname{prof}^*(\sh{F})$ instead of $\operatorname{prof}^*_X(\sh{F})$.

\oldpage[IV-3]{111}

\begin{remarks}[10.8.2]
\label{IV.10.8.2}
(i) At every closed point $x\in X$, the rectified depth equals the depth.

(ii) Let $Y$ be a closed sub-prescheme of $X$, $j:Y\to X$ the canonical injection.
We know \sref{IV.5.7.3},~(vi) that $\operatorname{prof}(\sh{G}_x)=\operatorname{prof}(j_*(\sh{G})_x)$ for every $x\in Y$; we therefore deduce that $\operatorname{prof}^*_x(\sh{G})=\operatorname{prof}^*_x(j_*(\sh{G}))$ for every $x\in Y$.

(iii) The notion of rectified depth is of interest only when it is of a \emph{local} character, i.e.\ when it does not change when one replaces $X$ by a neighbourhood open $U$ of $x$.
This obviously requires that $x$ not be isolated in $\overline{\{x\}}$ when $x$ is not a closed point of $X$, i.e.\ most often, that $X$ be a \emph{Jacobson prescheme} \sref{IV.10.4.5.1}; furthermore, it will also be necessary to know that $\dim(U)=\dim(X)$ for every dense open $U$ in $X$, and one must therefore also suppose that $X$ satisfies conditions $2^\circ$ and $3^\circ$ of \sref{IV.10.6.1}.
\end{remarks}

\begin{lemma}[10.8.3]
\label{IV.10.8.3}
Let $X$ be a regular and biequidimensional prescheme, $\sh{F}$ a coherent $\sh{O}_X$-Module.
Then, for every $x\in X$, we have
\begin{equation}
\label{IV.10.8.3.1}
\tag{10.8.3.1}
\operatorname{prof}^*_x(\sh{F})=\dim(X)-\dim.\operatorname{proj}(\sh{F}_x).
\end{equation}
\end{lemma}

\begin{proof}
Indeed, as $X$ is biequidimensional, we have \sref[0]{0.14.3.5.1}
\[
  \dim(\overline{\{x\}})=\dim(X)-\operatorname{codim}(\overline{\{x\}},X)=\dim(X)-\dim(\sh{O}_{X,x})
\]
by virtue of \sref{IV.5.1.2}.
On the other hand, as $X$ is regular, we have by \sref[0]{0.17.3.4}
\[
  \operatorname{prof}(\sh{F}_x)=\dim(\sh{O}_{X,x})-\dim.\operatorname{proj}(\sh{F}_x)
\]
whence the lemma.
\end{proof}

\begin{corollary}[10.8.4]
\label{IV.10.8.4}
Under the hypotheses of \sref{IV.10.8.3}, the function $x\mapsto\operatorname{prof}^*_x(\sh{F})$ is lower semi-continuous.
\end{corollary}

\begin{proof}
This follows from \sref{IV.10.8.3.1}, since $x\mapsto\dim.\operatorname{proj}(\sh{F}_x)$ is upper semi-continuous \sref{IV.6.11.1}.
\end{proof}

\begin{proposition}[10.8.5]
\label{IV.10.8.5}
Let $S$ be a prescheme locally Noetherian, $X$ a prescheme locally of finite type over $S$, $\sh{F}$ a coherent $\sh{O}_X$-Module.
Suppose that $S$ satisfies the conditions: $1^\circ$~$S$ is a Jacobson prescheme; $2^\circ$~$S$ is regular; $3^\circ$~the irreducible components of $S$ are equicodimensional.
Then the function $x\mapsto\operatorname{prof}^*_x(\sh{F})$ is semi-continuous inferiorly; in other terms, for every integer $n$, the set $U_n$ of $x\in X$ such that $\operatorname{prof}^*_x(\sh{F})\geq n$ is open.
\end{proposition}

\begin{proof}
As the local rings $\sh{O}_{S,s}$ of $S$ are regular, they are universally catenary \sref{IV.5.6.4}; in other words, $S$ satisfies conditions $1^\circ$, $2^\circ$, and $3^\circ$ of \sref{IV.10.6.1}, hence the same is true of $X$ \sref{IV.10.6.1},~(i).
The notion of rectified depth being then of local character \sref{IV.10.8.2},~(iii), one can restrict to the case where $S=\operatorname{Spec}(A)$ and $X=\operatorname{Spec}(B)$ are affine, $A$ being a regular ring and $B$ a quotient of $A$-algebra of finite type, hence a quotient of a polynomial ring $C=A[T_1,\ldots,T_n]$, and this last ring is regular \sref[0]{0.17.3.7}.
One can therefore suppose that $X$ is a closed sub-prescheme of a regular prescheme $Y$ also satisfying conditions $1^\circ$, $2^\circ$, and $3^\circ$ of \sref{IV.10.6.1}; taking into account the remark \sref{IV.10.8.2},~(ii), one is thus brought to the case where $X$ is in addition regular and Noetherian.
But as the local rings of $X$ are then integral, the irreducible components of $X$ are
\oldpage[IV-3]{112}
open \sref[I]{I.6.1.10}, and one can therefore also suppose $X$ irreducible.
Then, as the local rings of $X$ are catenary \sref[0]{0.16.5.12}, the hypothesis $3^\circ$ of \sref{IV.10.6.1} implies that $X$ is biequidimensional (\sref{IV.5.1.5} and \sref[0]{0.14.3.3}); it therefore suffices to apply \sref{IV.10.8.4}.
\end{proof}

One notes that if $S$ is the spectrum of a field or of $\mathbf{Z}$, it satisfies the conditions of \sref{IV.10.8.5}.

\begin{corollary}[10.8.6]
\label{IV.10.8.6}
The hypotheses being those of \sref{IV.10.8.5}, for every $x\in X$, the number $\operatorname{prof}^*_x(\sh{F})$ is the unique integer $n$ having the following property: there exists an open neighbourhood $U$ of $x$ such that for every point $x'\in U\cap\overline{\{x\}}$, closed in $U$, we have $\operatorname{prof}(\sh{F}_{x'})=n$.
In particular, for every $x$ such that $\operatorname{prof}^*_x(\sh{F})\geq m$ (resp.\ $\operatorname{prof}^*_x(\sh{F})\leq m$), it is necessary and sufficient that there exist an open $V$ of $X$ such that, for every point $x'\in V\cap\overline{\{x\}}$, closed in $V$, we have $\operatorname{prof}(\sh{F}_{x'})\geq m$ (resp.\ $\operatorname{prof}(\sh{F}_{x'})\leq m$).
\end{corollary}

\begin{proof}
Indeed, one can restrict to the case where $x$ is not a closed point; if $\operatorname{prof}^*_x(\sh{F})=n$, the set of $y\in X$ such that $\operatorname{prof}^*_y(\sh{F})\leq n$ is closed by virtue of \sref{IV.10.8.5}, hence contains $\overline{\{x\}}$, and by virtue of the lower semi-continuity of $\operatorname{prof}^*_y(\sh{F})$, there exists an open neighbourhood $U$ of $x$ such that $\operatorname{prof}^*_y(\sh{F})=n$ for every $y\in U\cap\overline{\{x\}}$, hence $\operatorname{prof}(\sh{F}_y)=n$ if $x'\in U\cap\overline{\{x\}}$ is closed (since the notion of rectified depth is local).
\end{proof}

For preschemes satisfying the hypotheses of \sref{IV.10.8.5}, the notion of rectified depth can therefore be defined by means of the values of the depth at the \emph{closed points} of $X$ (the latter forming a set which is very dense in every closed subset of $X$).

\begin{proposition}[10.8.7]
\label{IV.10.8.7}
Let $S$ be a prescheme locally Noetherian satisfying conditions $1^\circ$, $2^\circ$, and $3^\circ$ of \sref{IV.10.6.1}.
Let $X$ be a prescheme locally of finite type over $S$, $\sh{F}$ a coherent $\sh{O}_X$-Module, $Y=\operatorname{Supp}(\sh{F})$.
Then, for every $x\in Y$, we have
\begin{equation}
\label{IV.10.8.7.1}
\tag{10.8.7.1}
\operatorname{prof}^*_x(\sh{F})=\dim_x(Y)-\operatorname{coprof}(\sh{F}_x).
\end{equation}
\end{proposition}

\begin{proof}
Indeed, by definition, $\operatorname{coprof}(\sh{F}_x)=\dim(\sh{F}_x)-\operatorname{prof}(\sh{F}_x)$, and it follows from \sref{IV.10.6.4} that $\dim_x(Y)=\dim(\overline{\{x\}})+\dim(\sh{F}_x)$; whence \sref{IV.10.8.7.1} by definition of $\operatorname{prof}^*_x(\sh{F})$.
\end{proof}

\begin{corollary}[10.8.8]
\label{IV.10.8.8}
Under the hypotheses on $f$ and $\sh{F}$ being those of \sref{IV.9.9.1}, the function $x\mapsto\operatorname{prof}^*_x(\sh{F}_{f(x)})$ is constructible.
\end{corollary}

\begin{proof}
For every $s\in S$, the fibre $X_s$, being locally of finite type over $k(s)$, is a Jacobson prescheme \sref{IV.10.4.7}. Since $\operatorname{Spec}(k(s))$ satisfies the conditions of \sref{IV.10.6.1}, one has
\[
  \operatorname{prof}^*_x(\sh{F}_{f(x)})
  =\dim_x(\operatorname{Supp}(\sh{F}_{f(x)}))
   -\operatorname{coprof}((\sh{F}_{f(x)})_x)
\]
by \sref{IV.10.8.7}.
Now, if $Z=\operatorname{Supp}(\sh{F})$, we have $Z_{f(x)}=\operatorname{Supp}(\sh{F}_{f(x)})$ \sref[I]{I.9.1.13} and $Z$ is locally constructible \sref{IV.8.9.1}; therefore the functions $x\mapsto\dim_x(\operatorname{Supp}(\sh{F}_{f(x)}))$ and $x\mapsto\operatorname{coprof}((\sh{F}_{f(x)})_x)$ are locally constructible (\sref{IV.9.9.1} and \sref{IV.9.9.3}), which proves the proposition.
\end{proof}

\subsection{Maximal spectra and ultra-preschemes}
\label{subsection:IV.10.9}

The results of this section will not be used in what follows.

\begin{env}[10.9.1]
\label{IV.10.9.1}
Let $X$ be a Jacobson prescheme, and let $S(X)$ be the \emph{ringed space} whose underlying space is the subspace of closed points of $X$, and the sheaf of rings the sheaf induced on this subspace by $\sh{O}_X$, in other words the sheaf $\psi^*(\sh{O}_X)$, where $\psi:S(X)\to X$ denotes the canonical injection.
As $\psi$ is a quasi-
\oldpage[IV-3]{113}
homeomorphism, we have seen \sref{IV.10.2.8},~(ii) that if $\theta:\sh{O}_X\to\psi_*(\sh{O}_{S(X)})$ is the homomorphism of sheaves of rings such that $\theta^\sharp:\psi^*(\sh{O}_X)\to\sh{O}_{S(X)}$ is the identity, then $j_x=(\psi,\theta)$ is a \emph{quasi-isomorphism} of ringed spaces, and $\sh{F}\mapsto j_X^*(\sh{F})$ is an equivalence of the category of $\sh{O}_X$-Modules and that of $\sh{O}_{S(X)}$-Modules.
It is clear that under this equivalence, locally free (resp.\ coherent) $\sh{O}_X$-modules correspond to locally free (resp.\ coherent) $\sh{O}_{S(X)}$-modules. Furthermore, if $\sh{L}$ is a locally free $\sh{O}_X$-module and $U$ is an open subset of $X$ such that $\sh{L}|_U\simeq(\sh{O}_X|_U)^n$, then
\[
  j_X^*(\sh{L})|_{U\cap S(X)}
  \simeq \bigl(\sh{O}_{S(X)}|_{U\cap S(X)}\bigr)^n.
\]
\end{env}

\begin{env}[10.9.2]
\label{IV.10.9.2}
Let $X$, $Y$ be two Jacobson preschemes and $f=(p,\lambda):X\to Y$ a morphism \emph{locally of finite type}.
One has seen \sref{IV.10.4.7} that $p(S(X))\subset S(Y)$ and by restriction of $p$ to $S(X)$, one defines therefore a continuous map $S(p):S(X)\to S(Y)$.
For every open $V$ of $Y$, one defines by composition a homomorphism of rings
\[
  \Gamma(V\cap S(Y),\sh{O}_{S(Y)})\to\Gamma(V,\sh{O}_Y)\xrightarrow{\lambda_V}\Gamma(f^{-1}(V),\sh{O}_X)\to\Gamma(f^{-1}(V)\cap S(X),\sh{O}_{S(X)})
\]
where the two extreme isomorphisms have been defined in \sref{IV.10.9.1}; it is clear that this defines a homomorphism of sheaves of rings $S(\lambda):S(\sh{O}_{S(Y)})\to\rho_*(\sh{O}_{S(X)})$ (recalling that the open subsets of $S(X)$ (resp.\ $S(Y)$) correspond bijectively to those of $X$ (resp.\ $Y$) \sref{IV.10.2.1}); one obtains thus a morphism of ringed spaces $S(f)=(S(p),S(\lambda)):(S(X),\sh{O}_{S(X)})\to(S(Y),\sh{O}_{S(Y)})$ such that the diagram
\[
  \xymatrix{
    S(X)\ar[r]^{S(f)}\ar[d]_{j_X} & S(Y)\ar[d]^{j_Y}\\
    X\ar[r]^{f} & Y
  }
\]
is commutative; furthermore, if $Z$ is a third Jacobson prescheme and $g:Y\to Z$ a morphism locally of finite type, $S(g\circ f)=S(g)\circ S(f)$.
One has thus defined a \emph{covariant functor} $S:\mathbf{C}\to\mathbf{C}'$, where $\mathbf{C}$ is the category of \emph{ringed spaces in local rings}, and $\mathbf{C}'$ the category whose objects are the \emph{Jacobson preschemes} and the morphisms are the morphisms locally of finite type between Jacobson preschemes.
\end{env}

\begin{env}[10.9.3]
\label{IV.10.9.3}
Let us now determine the subcategory $\mathbf{C}''$ of $\mathbf{C}'$ formed by the ringed spaces isomorphic to $S(X)$.
Suppose first that $X=\operatorname{Spec}(A)$, where $A$ is a Jacobson ring; then $S(X)$ is the set of \emph{maximal} ideals of $A$, equipped: $1^\circ$~with the topology induced by that of $X$, so that a base for this topology is formed by the $D^m(h)=D(h)\cap S(X)$, i.e.\ the set of maximal ideals $\mathfrak{m}$ of $A$ such that $h\notin\mathfrak{m}$; $2^\circ$~with the sheaf of rings $\sh{O}_{S(X)}$ such that $\Gamma(D^m(h),\sh{O}_{S(X)})=A_h$.
We will say that this ringed space is the \emph{maximal spectrum} of the Jacobson ring $A$ and denote it $\operatorname{Spm}(A)$.

One notes that $j:D(h)\to X$ is the canonical injection, the ringed space induced on $D^m(h)$ by $S(X)$ and $S(D(h))$ are canonically isomorphic, and that the canonical injection $D^m(h)\to S(X)$ of ringed spaces is equal to $S(j)$.

Let $B$ be a second Jacobson ring, $Y=\operatorname{Spec}(B)$, $\varphi:B\to A$ a ring homomorphism making $A$ a $B$-algebra of \emph{finite type}, $f=({}^a\varphi,\widetilde{\varphi}):X\to Y$ the corresponding morphism of preschemes, and
\[
  S(f):\operatorname{Spm}(A)\to\operatorname{Spm}(B)
\]
the morphism of ringed spaces corresponding to $f$.
It is clear that $S(f)=(\psi,\theta)$ is a morphism of ringed spaces in local rings, i.e.\ $(\operatorname{Err}_{11},1.8.2)$ for every $x\in\operatorname{Spm}(A)$, $\theta^\sharp_x$ is a \emph{local} homomorphism.
Conversely:
\end{env}

\begin{proposition}[10.9.4]
\label{IV.10.9.4}
Let $A$, $B$ be two Jacobson rings.
If $u=(\psi,\theta):\operatorname{Spm}(A)\to\operatorname{Spm}(B)$ is a morphism of ringed spaces in local rings such that $\Gamma(\theta):\operatorname{B}\to A$ makes $A$ a $B$-algebra of finite type, there exists a morphism of preschemes $f:\operatorname{Spec}(A)\to\operatorname{Spec}(B)$ and a unique one such that $u=S(f)$.
\end{proposition}

\begin{proof}
The uniqueness of $f$ is evident, since if $f=({}^a\varphi,\widetilde{\varphi})$, one must have $\varphi=\Gamma(\theta)$; it remains to show that $S(f)$ is defined and that $u=S(f)$.
Now, the first assertion follows from the fact that $\varphi$ is supposed to make $A$ a $B$-algebra of finite type, the fact that $\theta^\sharp_x$ is a local homomorphism for every $x\in\operatorname{Spm}(A)$ allows one to repeat the argument of \sref[I]{I.1.7.3}, restricting to points $x$ of $\operatorname{Spm}(A)$, to show thus successively that $\psi(x)={}^a\varphi(x)$ for every $x\in\operatorname{Spm}(A)$, then that $\theta^\sharp_x=\varphi_x:B_{(x)}\to A_x$, which completes proving that $u=S(f)$.
\end{proof}

\oldpage[IV-3]{114}

\begin{env}[10.9.5]
\label{IV.10.9.5}
Consider now a ringed space $(X,\sh{O}_X)$; we will say that an open subset $U$ of $X$ is \emph{ultra-affine} if the ringed space induced $(U,\sh{O}_X|U)$ is isomorphic to a maximal spectrum $\operatorname{Spm}(A)$, where $A$ is a Jacobson ring.
We will say that $X$ is an \emph{ultra-prescheme} if every point of $X$ admits an ultra-affine open neighbourhood.
One shows, as in \sref[I]{I.2.1.3} and \sref{IV.2.1.4}, that the ultra-affine open subsets form a \emph{base} for the topology of $X$ and that $X$ is a Kolmogoroff space.
If $Y$ is a second ultra-prescheme, we will say that a morphism of ringed spaces $f:X\to Y$ is a \emph{morphism of ultra-preschemes} if it satisfies the following conditions: $1^\circ$~$f$ is a morphism of ringed spaces in local rings; $2^\circ$~for every $x\in X$, there is an ultra-affine open neighbourhood $V$ of $f(x)$ in $Y$ and an ultra-affine open $U$ of $x$ in $X$ such that $f(U)\subset V$ and the homomorphism $\Gamma(V,\sh{O}_Y)\to\Gamma(U,\sh{O}_X)$ corresponding to $f$ makes $\Gamma(U,\sh{O}_X)$ a $\Gamma(V,\sh{O}_Y)$-\emph{algebra of finite type}.

It is immediate that one defines well the composite of two such morphisms, being a third such by the final remark of \sref{IV.10.9.3}.
It is clear that the category $C''_0$ thus defined is a subcategory of $C'$ that contains $C''$; we now propose to show that $C''=C''_0$, in other words:
\end{env}

\begin{proposition}[10.9.6]
\label{IV.10.9.6}
The functor $X\mapsto S(X)$ of $\mathbf{C}$ into $\mathbf{C}''_0$ is an equivalence of categories.
\end{proposition}

\begin{proof}
$1^\circ$ Let us first show that the functor $X\mapsto S(X)$ is \emph{fully faithful}, in other words that for $X$, $Y$ in $\mathbf{C}$, the canonical map
\[
  \operatorname{Hom}_{\mathbf{C}}(X,Y)\to\operatorname{Hom}_{\mathbf{C}''_0}(S(X),S(Y))
\]
is \emph{bijective}.
In the first place it is \emph{injective}: let $f$, $g$ be two morphisms locally of finite type of $X$ into $Y$, and suppose that $S(f)=S(g)$; we have $f^{-1}(V)\cap S(X)=g^{-1}(V)\cap S(X)$ for every open $V$ of $Y$, hence \sref{IV.10.3.1} and \sref{IV.10.2.7} $f^{-1}(V)=g^{-1}(V)$; it suffices therefore to prove that for every affine open $V$ of $Y$, $f$ and $g$ coincide in $f^{-1}(V)=g^{-1}(V)$; in other words, one reduces to the case where $Y=\operatorname{Spec}(B)$ is the spectrum of a Jacobson ring $B$.
It suffices \sref[I]{I.2.2.4} to show that the homomorphisms of rings $B\to\Gamma(X,\sh{O}_X)$ corresponding to $f$ and $g$ are the same.
Now, for every $s\in B$, the images of $s$ under these homomorphisms are two sections of $\sh{O}_X$ over $X$ which, by hypothesis, induce the \emph{same} section over $S(X)$; we know \sref{IV.10.2.8} that these sections are identical, whence our assertion.
In the second place let us prove that every morphism $h:S(X)\to S(Y)$ (for the category $\mathbf{C}''_0$) is of the form $S(f)$, where $f:X\to Y$ is a morphism locally of finite type.
By hypothesis, there exists an ultra-affine open covering $(U')$ (resp.\ $(V'_\mu)$) of $S(X)$ (resp.\ $S(Y)$) such that for every $\lambda$, $h(U'_\lambda)$ is contained in one of the $V'_\mu$ and that it corresponds to the morphism $h_\lambda:U'_\lambda\to V'_\mu$ making the second ring an algebra of finite type over the first.
For every couple of indices $\alpha$, $\beta$, consider the open subset $U_{\alpha\beta}$ of $U_\alpha$ whose trace on $S(X)$ is $U'_\alpha\cap U'_\beta$; by virtue of $1^\circ$, the identical self-morphism of $U'_\alpha\cap U'_\beta$ is of the form $S(\theta_{\alpha\beta})$, where $\theta_{\alpha\beta}:U_{\alpha\beta}\to U_{\beta\alpha}$ is an isomorphism of preschemes.
We see immediately (by virtue of $1^\circ$) that the family $(\theta_{\alpha\beta})$ satisfies the glueing condition of recollement \sref[0]{0.4.1.7}, and that this family defines a prescheme $X$ such that we have $X'=S(X)$; it is clear then that the $U_\alpha$ are identified with affine open subsets of $X$, such that the $U'_\alpha$ are the restrictions of $S(X)$ to the $U_\alpha$.

$2^\circ$ It remains to prove that every ultra-prescheme $X'$ is of the form $S(X)$ for a Jacobson prescheme $X$ (which will necessarily be unique to isomorphism near, by virtue of $1^\circ$).
There are ultra-affine open subsets $(U'_\alpha)$ of $X'$, each of which is of the form $S(U_\alpha)$, $U_\alpha$ being the spectrum of a Jacobson ring.
For every couple of indices $\alpha$, $\beta$, consider the open $U_{\alpha\beta}$ of $U_\alpha$ whose trace $U'_\alpha$ is $U'_\alpha\cap U'_\beta$; by virtue of $1^\circ$, the identical self-morphism of $U'_\alpha\cap U'_\beta$ is of the form $S(\theta_{\alpha\beta})$, where $\theta_{\alpha\beta}:U_{\alpha\beta}\to U_{\beta\alpha}$ is an isomorphism of preschemes.
One easily verifies (by virtue of $1^\circ$) that the family $(\theta_{\alpha\beta})$ satisfies the glueing condition \sref[0]{0.4.1.7}, and that this family defines a prescheme $X$; it is clear then that we have $X'=S(X)$, which completes the proof.
\end{proof}

\subsection{Algebraic spaces of Serre}
\label{subsection:IV.10.10}

\begin{env}[10.10.1]
\label{IV.10.10.1}
The language introduced by Serre in (FAC) is sometimes convenient, notably in questions where the main interest is attached to the rational points over the base field $k$ (algebraically closed by hypothesis) of the ``algebraic varieties'' over $k$ that one considers.
We will sketch here this language by linking it to the considerations that precede, to allow the reader to translate the statements of Serre in the language of schemes.
It is moreover possible to develop the language of Serre also for preschemes over a non-algebraically closed field (and even over an Artinian ring); but this introduces considerable technical complications, and moreover, on a field of arbitrary base, the advantages (above all psychological ones) of the point of view of Serre disappear; therefore we will confine ourselves to the framework fixed by Serre.
The present section, just like the preceding one, will not be used in the sequel of this Treatise, and we will therefore content ourselves with brief indications.
\end{env}

\oldpage[IV-3]{115}

\begin{env}[10.10.2]
\label{IV.10.10.2}
Given an ultra-prescheme $R$, one can naturally (as in the whole category) define the notion of $R$-\emph{ultra-prescheme}.
Consider in particular an algebraically closed field $k$; $\operatorname{Spm}(k)$ is then identical to $\operatorname{Spec}(k)$; we will say that a $\operatorname{Spm}(k)$-ultra-prescheme $X'$ is a \emph{prealgebraic space over $k$}: it is a ringed space in local rings whose every point has an open neighbourhood isomorphic to the maximal spectrum of a $k$-algebra of finite type; it follows that it is the same as saying (by \sref{IV.10.9.6}) that $X'=S(X)$, where $X$ is a $k$-prescheme locally of finite type.
If $X$, $Y$, $Z$ are three $k$-preschemes locally of finite type, it is the same as to say of $X\times_k Y$ \sref[I]{I.1.3.4}, hence the same is true of $X'\times_k Y'$, in the category of prealgebraic spaces over $k$ (and also of the ``fibre product'' $X'\times_{k'}Y'$, where $X'$, $Y'$, $Z'$ are three prealgebraic spaces over $k$).
One can therefore define the diagonal morphism $\Delta_{X'}:X'\to X'\times_k X'$ (which is none other than $S(\Delta_X)$); one says that $X'$ is an \emph{algebraic space over $k$} if $X$ is a scheme, and it amounts to the same as saying that the image of $\Delta_{X'}$ is a \emph{closed} subset of $X'\times_k X'$.
\end{env}

\begin{env}[10.10.3]
\label{IV.10.10.3}
The simplifications that come from the hypothesis that $k$ is algebraically closed amount first of all to this: for every $k$-prescheme $X$ locally of finite type, there is a bijective correspondence between \emph{closed points} of $X_k$ (i.e.\ the \emph{points of $X$ with values in $k$} \sref[I]{I.3.4.4}) and \emph{rational points of $X$ over $k$} \sref[I]{I.3.4.5}, by virtue of \sref[I]{I.6.4.2}.
This shows in particular \sref[I]{I.3.4.3.1} that for two $k$-prealgebraic spaces $X'$, $Y'$, the underlying set product $X'\times Y'$ \emph{is} identical to the underlying set of the product of the underlying topological spaces of $X'$ and $Y'$ (but it is well understood that the topology of the latter is not the \emph{product} topology of the topologies of the underlying spaces, in general it is strictly finer than the latter).

On the other hand, the local rings $\sh{O}_x$ at the points of a prealgebraic space $X'$ over $k$ are $k$-algebras whose residue field is \emph{isomorphic to $k$}; if $A$ and $B$ are two such $k$-algebras, a $k$-homomorphism $\varphi:A\to B$ is necessarily \emph{local}: indeed, if an element $x$ of the maximal ideal $A$ were such that $\varphi(x)$ is invertible, there would exist $\lambda\in k$ nonzero such that $\varphi(x-\lambda.1)$ belongs to the maximal ideal of $B$, which is absurd since $x-\lambda.1$ is invertible.
We conclude immediately from this that if $X'=S(X)$, $Y'=S(Y)$ are two prealgebraic spaces over $k$, every morphism $X'\to Y'$ of \emph{$k$-ringed spaces} is also a morphism of ringed spaces \emph{in local rings}, hence every morphism of $k$-preschemes $X'\to Y'$ of $k$-ringed spaces is also, in the notations above, of the form $S(f)$, where $f:X\to Y$ is a morphism of $k$-preschemes.

Finally, for every open $U$ of $X'$, and every $x\in U$, $s(x)$ \sref[0]{0.5.5.1} is identified with an element of $U$ in $k$, in other words a map $\overline{s}:x\to s(x)$ of $U$ in $k$; as the map $h_U:s\to\overline{s}$ is evidently a homomorphism of rings $\Gamma(U,\sh{O}_{X'})\to\Gamma(U,\sh{A}(X'))$ (the \emph{germs of maps} of $X'$ \emph{in $k$}) and commutes with restrictions to a neighbourhood $V\subset U$, the $h_U$ define a homomorphism of \emph{sheaves} of rings $h:\sh{O}_{X'}\to\sh{A}(X')$.
If one takes $U$ to be ultra-affine, $\operatorname{Spm}(A)$, where $A$ is a Jacobson ring, to say that $\overline{s}=0$ means that for every maximal ideal $\mathfrak{m}$ of $A$, $s$ belongs to $\mathfrak{m}$, or equivalently that $s$ is in the \emph{radical} of $A$; but as $A$ is a Jacobson ring, its radical is its nilradical, hence $\overline{s}=0$ if and only if $s$ is nilpotent.
\end{env}

\begin{env}[10.10.4]
\label{IV.10.10.4}
One says that the $k$-prealgebraic space $X'=S(X)$ is \emph{reduced} if it is reduced at each of its points, i.e.\ the set of points $x\in X$ where $X$ is reduced is open \sref[0]{0.5.2.2}, its complement contains at least one closed point $x'$ of $X$; as each of the local rings $\sh{O}_{x'}$ is reduced, one sees in \sref{IV.10.10.3} that for the homomorphism $h:\sh{O}_{X'}\to\sh{A}(X')$ to be \emph{injective}, it is necessary and sufficient that $X'$ be \emph{reduced}.
In (FAC), Serre restricts himself in fact to prealgebraic \emph{reduced} spaces, which allows him to define $\sh{O}_{X'}$ as a \emph{sub-sheaf} of $\sh{A}(X')$.
One notes that if $X'$ and $Y'$ are two $k$-spaces prealgebraic reduced, it is the same of $X'\times_k Y'$: indeed, this amounts to saying that if $A$ and $B$ are two reduced $k$-algebras of finite type, $A\otimes_k B$ is also reduced; but we have just seen that the radicals of $A$ and $B$ are $0$, and as $k$ is algebraically closed, $A$ and $B$ are ``separable'' algebras over $k$ in the sense of Bourbaki (\emph{Alg.}, chap.~VIII, \S\,7, n$^\circ$\,5, prop.~5); hence $A\otimes_k B$ is without radical (\emph{loc.\ cit.}, n$^\circ$\,6, cor.~3 of th.~3), and since it is a Jacobson ring, it is reduced.
However, if $Z'$ is a third prealgebraic space reduced over $k$, the ``fibre product'' $X'\times_{Z'}Y'$ of two prealgebraic reduced spaces over $Z'$ is not in general reduced, which implies that the category of these spaces is insufficient in many questions (notably in the theory of algebraic groups).
But as we have seen above, one can keep the language of Serre without restricting oneself, as this latter (who in addition considers only quasi-compact \emph{prealgebraic} spaces), to the case of reduced prealgebraic spaces.
\end{env}

\oldpage[IV-3]{116}

\begin{env}[10.10.5]
\label{IV.10.10.5}
Finally, one can also consider ultra-preschemes over an arbitrary field $k$ while
retaining a language close to Serre's, and introducing, as in Weil, a fixed
algebraically closed extension $K$ of $k$ (chosen sufficiently large---for example,
of infinite transcendence degree over $k$---so as to have enough ``generic points''
in Weil's sense).
To every prescheme $X$ locally of finite type over $k$ one then associates the set
\[
 S_K(X)=S(X\otimes_k K)
\]
of points of $X$ with values in $K$.
There is a canonical map $j:S_K(X)\to X$, which is a quasi-homeomorphism when
$S_K(X)$ is given the inverse-image topology from $X$ under $j$; endowing $S_K(X)$
with the sheaf of $k$-algebras $j^*(\sh{O}_X)$ gives a subcategory of the category
of $k$-ringed spaces in local rings, which one might call the category of
\emph{$(k,K)$-prealgebraic spaces}.
Products can still be defined there, and the results of \sref{IV.10.10.3} and
\sref{IV.10.10.4} can be generalized (with $\sh{A}(X')$ replaced here by the sheaf
$\sh{A}_K(X')$ of germs of maps from $X'$ into $K$).
This point of view, however, assigns an artificial role to an arbitrarily chosen
overfield of $k$, and we mention it only to reject it.
\end{env}
